% Options for packages loaded elsewhere
\PassOptionsToPackage{unicode}{hyperref}
\PassOptionsToPackage{hyphens}{url}
\PassOptionsToPackage{dvipsnames,svgnames,x11names}{xcolor}
\documentclass[
  a4paper,
]{article}
\usepackage{xcolor}
\usepackage[margin=25mm]{geometry}
\usepackage{amsmath,amssymb}
\setcounter{secnumdepth}{-\maxdimen} % remove section numbering
\usepackage{iftex}
\ifPDFTeX
  \usepackage[T1]{fontenc}
  \usepackage[utf8]{inputenc}
  \usepackage{textcomp} % provide euro and other symbols
\else % if luatex or xetex
  \usepackage{unicode-math} % this also loads fontspec
  \defaultfontfeatures{Scale=MatchLowercase}
  \defaultfontfeatures[\rmfamily]{Ligatures=TeX,Scale=1}
\fi
\usepackage{lmodern}
\ifPDFTeX\else
  % xetex/luatex font selection
  \setmainfont[]{DejaVu Serif}
  \setsansfont[]{DejaVu Sans}
  \setmonofont[]{DejaVu Sans Mono}
\fi
% Use upquote if available, for straight quotes in verbatim environments
\IfFileExists{upquote.sty}{\usepackage{upquote}}{}
\IfFileExists{microtype.sty}{% use microtype if available
  \usepackage[]{microtype}
  \UseMicrotypeSet[protrusion]{basicmath} % disable protrusion for tt fonts
}{}
\makeatletter
\@ifundefined{KOMAClassName}{% if non-KOMA class
  \IfFileExists{parskip.sty}{%
    \usepackage{parskip}
  }{% else
    \setlength{\parindent}{0pt}
    \setlength{\parskip}{6pt plus 2pt minus 1pt}}
}{% if KOMA class
  \KOMAoptions{parskip=half}}
\makeatother
\usepackage{longtable,booktabs,array}
\newcounter{none} % for unnumbered tables
\usepackage{calc} % for calculating minipage widths
% Correct order of tables after \paragraph or \subparagraph
\usepackage{etoolbox}
\makeatletter
\patchcmd\longtable{\par}{\if@noskipsec\mbox{}\fi\par}{}{}
\makeatother
% Allow footnotes in longtable head/foot
\IfFileExists{footnotehyper.sty}{\usepackage{footnotehyper}}{\usepackage{footnote}}
\makesavenoteenv{longtable}
\usepackage{graphicx}
\makeatletter
\newsavebox\pandoc@box
\newcommand*\pandocbounded[1]{% scales image to fit in text height/width
  \sbox\pandoc@box{#1}%
  \Gscale@div\@tempa{\textheight}{\dimexpr\ht\pandoc@box+\dp\pandoc@box\relax}%
  \Gscale@div\@tempb{\linewidth}{\wd\pandoc@box}%
  \ifdim\@tempb\p@<\@tempa\p@\let\@tempa\@tempb\fi% select the smaller of both
  \ifdim\@tempa\p@<\p@\scalebox{\@tempa}{\usebox\pandoc@box}%
  \else\usebox{\pandoc@box}%
  \fi%
}
% Set default figure placement to htbp
\def\fps@figure{htbp}
\makeatother
\setlength{\emergencystretch}{3em} % prevent overfull lines
\providecommand{\tightlist}{%
  \setlength{\itemsep}{0pt}\setlength{\parskip}{0pt}}
\usepackage{luatexja}
\usepackage{bookmark}
\IfFileExists{xurl.sty}{\usepackage{xurl}}{} % add URL line breaks if available
\urlstyle{same}
\hypersetup{
  pdftitle={The Eleventh Thought Experiment: Emission Records of Two-Body Systems without a Background Spacetime},
  pdfauthor={Noriaki Kihara (WF System Co., Ltd.) ORCID 0009-0004-6753-4020},
  colorlinks=true,
  linkcolor={Maroon},
  filecolor={Maroon},
  citecolor={Blue},
  urlcolor={Blue},
  pdfcreator={LaTeX via pandoc}}

\title{The Eleventh Thought Experiment: Emission Records of Two-Body
Systems without a Background Spacetime}
\usepackage{etoolbox}
\makeatletter
\providecommand{\subtitle}[1]{% add subtitle to \maketitle
  \apptocmd{\@title}{\par {\large #1 \par}}{}{}
}
\makeatother
\subtitle{- Computing Standard Theory (Dirac--Coulomb Bound States, 1PN
Two-Body Gravity, Einstein A Coefficients) Referenced Only to the Rest
Frame of an Absorber at Infinity, and Reading Out Emission--Absorption
Relations for the Hydrogen Atom, Electron--Electron, Proton--Proton and
the Hydrogen Molecule -}
\author{Noriaki Kihara (WF System Co., Ltd.) ORCID 0009-0004-6753-4020}
\date{October 7, 2026 Version DOI 10.5281/zenodo.23199300}

\ltjsetparameter{jacharrange={-2,-3}}
\begin{document}
\maketitle

\textbf{Author:} Noriaki Kihara\\
\textbf{ORCID:} 0009-0004-6753-4020\\
\textbf{Version DOI:} 10.5281/zenodo.23199300\\
\textbf{Concept DOI:} 10.5281/zenodo.23199299\\
\textbf{Version:} v1.0\\
\textbf{Date:} 2026-10-07\\
\textbf{Related papers:} the design document (Paper 0) {[}S1{]}, the
sixth thought experiment v1.1 {[}S2{]}, its supplement {[}S3{]}, the
composition of central projection {[}S4{]}, the radial projection note
{[}S5{]} (all as references; this paper is a separate line)

\textbf{Keywords:} two-body system, emission record, absorber, Einstein
A coefficients, master equation, Dirac--Coulomb, recoil corrections, 1PN
gravity, gravitational waves, hydrogen atom, hydrogen molecule, Mott
scattering, radiative association, cosmic microwave background,
reproducibility

\begin{center}\rule{0.5\linewidth}{0.5pt}\end{center}

\section{Abstract}\label{abstract}

A re-examination of the previous paper (the sixth thought experiment
{[}S2{]} and its supplement {[}S3{]}) showed that gravity, the Coulomb
force and the gravitational-wave and electromagnetic radiation were not
updated there by the distribution-array operation of the state
interaction, but by known equations supplied from outside. The seventh
to tenth thought experiments therefore tried to derive the dynamics of
gravity and the Coulomb force from the distribution-array interaction,
and none of them reached a two-body orbit with gravity included.
Changing course, the known Einstein--Maxwell equations were expanded at
initialization and implemented as discrete variational maps to be
replaced by distribution-array operations; the run with an absorber did
not reach a lock and collapsed, and the goal was not met.

This paper is a separate line. Setting aside the derivation of the
dynamics by the distribution-array operation of a universal interaction,
it builds a program that computes the known standard theory - the
Dirac--Coulomb bound states of the electron with recoil and QED-type
corrections, the 1PN two-body gravity of general relativity with
quadrupole gravitational waves, and the master equation of Einstein A
coefficients - without holding any coordinate grid or metric of a
background spacetime in the state. The units are \(G=c=q_0=1\) (\(q_0\)
is one third of the elementary charge); the inputs are the
fine-structure constant \(\alpha\), the mass ratio \(\mu_p=m_p/m_e\),
the ratio \(\mu=m_e\sqrt{G}/q_0\) that carries the strength of gravity,
and the proton \(g\)-factor, all at their known values. The only
reference frame is the rest frame of an absorber at infinity, which
corresponds to the asymptotic region of general relativity. Readouts use
only the record of emission and absorption (frequency, rate, branching,
time differences, single or pair, direction, polarization); no internal
state is read.

Accordingly this paper claims no new discovery. It is the record of an
experiment that made known knowledge computable under restricted
conditions, and at the same time it supplies the reference (the answer
table) of emission records that a future distribution-array operation
must reproduce.

The experiments were run on four systems: the cascade of the hydrogen
atom (proton and electron) from 5p; the same-sign pairs
electron--electron and proton--proton (no bound levels; only the in--out
quantities, the Mott cross section and quadrupole bremsstrahlung); and
hydrogen atom--hydrogen atom (the rovibrational cascade on the singlet
X¹Σg⁺ and radiative association from the continuum; the potential and
the quadrupole moment are literature values). The absorber is placed at
infinity, far enough not to affect the pair, as a complete absorber
covering all directions, and the isotropic case (\(T=0\) and 2.7255 K)
is compared with an anisotropy close to the actual universe (a dipole
from the centre-of-mass velocity 370 km/s, a quadrupole
\(a_2=4\times10^{-6}\), \(N=10^{60}\) cells).

Main numbers. Hydrogen atom: 50 levels (\(n\le5\)), total emission from
5p₃/₂ F=1 13.054539268 eV (equal to the level drop, difference
\(4\times10^{-15}\)), of which E1 11.84, two-photon 1.21 (11.9 \% of
cascades pass through 2s), M1 \(2.4\times10^{-6}\), E2
\(1.5\times10^{-6}\), gravitational waves \(1.4\times10^{-48}\) eV.
Comparison with external values: a 0.12 MHz mismatch in the 1S
ionization energy (the remainder after removing QED and nuclear size),
all digits of the recoil terms in the table of Eides et al.~reproduced,
E1 A coefficients within \(1.2\times10^{-4}\) of NIST, the Lyman-α width
99.71 MHz (measured 99.7). With the absorber at 2.7255 K the 1s
hyperfine doublet equilibrates to 0.745286 : 0.254714
\(=3e^{-\hbar\omega/kT}\) (relaxation \(6.9\times10^4\) years); the
centre-of-mass velocity 370 km/s is recovered to 0.09 m/s from the
Doppler record of 200 atoms, and from 29 events of the 21 cm line of a
single atom, together with the history of its kicks, to
\(2.4\times10^{-10}\) m/s. The steady alignment of 1s F=1 due to the
anisotropy of the background is \(5\times10^{-9}\) (dipole only) to
\(2.5\times10^{-8}\) (with the quadrupole). Electron--electron and
proton--proton: the net potential is positive, there are no bound
levels, and the only things left in the record are the Mott cross
section (one half of the distinguishable-particle value at 90°) and
quadrupole bremsstrahlung (\(5\times10^{-16}\) / \(5\times10^{-23}\)
photons per collision). Hydrogen molecule: 301 levels from a sinc-DVR on
the BO + adiabatic potential, \(D_0=36118.364\) cm⁻¹ (fully corrected
value 36118.070), E2 A coefficients agreeing with the 4669 lines of
Roueff et al.~(2019) to a median \(8\times10^{-4}\). Since 2.7255 K is
1/310 of the lowest transition, the background excites nothing, and the
velocity recovery from the Doppler record reaches 0.002 m/s because the
line widths are \(10^{-22}\). Radiative association H + H → H₂ + γ has a
capture probability \(10^{-19}\) per pass and
\(k(T)=2\)--\(4\times10^{-29}\) cm³ s⁻¹. Inter-atomic gravity,
\(10^{-31}\) cm⁻¹ per level, is kept as a separate account.

The absorber is not a sphere but the family of cuts of each emission
event's future light cone at infinity. Cut by the time slice of the
absorber's rest frame it is a sphere; cut by the time slice of the
two-body rest frame it is the ellipsoid
\(r(\theta)=t_0/(1-\beta\cos\theta)\), and the Doppler shift and
aberration of layer 1 are the geometry of this cut. Since the phase is
invariant along null geodesics, the record keeps the phase relations at
emission. ``No background spacetime'' and ``a complete absorber at
infinity'' do not contradict each other: the latter is a map onto the
null surface of the light cone, the same operation as the central and
radial projections of this research series.

\begin{center}\rule{0.5\linewidth}{0.5pt}\end{center}

\section{0. Position of this paper}\label{position-of-this-paper}

Under the design document {[}S1{]}, this research series has aimed to
close the motion of two bodies into one map that generates the next
state from the state alone. The sixth thought experiment {[}S2{]} put
gravity and the Coulomb force, gravitational-wave and electric-dipole
radiation into a single transition with eight persistent states, but its
update rule for radiation was the known multipole formula supplied from
outside. The attempts to derive the dynamics including radiation from
the distribution-array operation itself (thought experiments 7--10,
unpublished) and to put the Einstein--Maxwell equations into discrete
variational maps expanded at initialization (same, unpublished) both
failed to reach the intended result.

This paper does not aim at that derivation. Instead it fixes, as a
program that holds no background spacetime in its state, the emission
record that standard theory gives for the same two-body systems. If a
future distribution-array operation is correct, it must reproduce this
record. This paper supplies the target of that comparison.

By ``standard theory'' this paper means the Dirac--Coulomb bound states
of the electron and the two-body recoil corrections {[}4, 6, 8{]}, the
higher orders of retardation including the Bethe logarithm {[}3,
11--19{]}, the 1PN two-body Hamiltonian of general relativity {[}21,
22{]} and quadrupole radiation {[}27{]}, Einstein A coefficients with
the detailed balance of stimulated emission and absorption {[}28, 29{]},
and for the hydrogen molecule the Born--Oppenheimer potential with the
adiabatic correction {[}46--48{]} and the quadrupole moment {[}49{]}.
The field equations themselves are not evolved in time.

\begin{center}\rule{0.5\linewidth}{0.5pt}\end{center}

\section{1. Motivation and purpose}\label{motivation-and-purpose}

\subsection{1.1 Re-examination of the previous
paper}\label{re-examination-of-the-previous-paper}

The \texttt{transition(z)} of the sixth thought experiment {[}S2{]}
produces the next state from the eight persistent states
\(Z_8=(U,P,E,H,Q,N,C,D)\). The decrease of \(E\) and \(H\) by radiation,
however, was given as rates from outside by the known GW-quadrupole and
EM-dipole radiation formulas, not produced by a distribution-array
operation of the state interaction. The conservative part of gravity and
the Coulomb force was likewise a known effective potential turned into
coefficients at initialization. Thus what was ``closed in the state
alone'' was the form of the bookkeeping, not a derivation of the
interaction.

\subsection{1.2 Outcome of thought experiments 7--10 and of the
Einstein--Maxwell discrete
variation}\label{outcome-of-thought-experiments-710-and-of-the-einsteinmaxwell-discrete-variation}

The seventh (self-interaction aa, bb), eighth (cross coupling of aa, ab,
bb), ninth (finite-history readout by harmonic interference) and tenth
(finite history of the Liénard--Wiechert field and the Lorentz--Dirac
self-force, Regge--Wheeler--Zerilli, the coupled Reissner--Nordström
problem, exact anchors such as Kastor--Traschen) all failed to reach a
two-body orbit with gravity included. Next the Einstein--Maxwell
equations were expanded around the initial configuration (the EIH /
Infeld--Wallace procedure), implemented in a discrete action with a
variable step, and replaced by distribution-array operations; the run
with an absorber (200,000 periods) did not reach a phase lock and headed
for a Rutherford-type collapse. These records are kept in the repository
(§8) and are not cited in this paper.

\subsection{1.3 Purpose of this paper}\label{purpose-of-this-paper}

To turn the known standard theory into a numerically computable program
under the following conditions.

\begin{enumerate}
\def\labelenumi{\arabic{enumi}.}
\tightlist
\item
  No coordinate grid or metric of a background spacetime is held in the
  state. The only reference frame is the rest frame of an absorber at
  infinity (§2.2).
\item
  The units are \(G=c=q_0=1\); the inputs are \(\alpha\), \(\mu_p\),
  \(\mu\), \(g_p\) and, for H₂, literature potential curves (§2.1).
\item
  Readouts are made only from the record of emission and absorption
  (§2.3).
\item
  The absorber is given as isotropic (\(T=0\), 2.7255 K) and as
  anisotropic (dipole, quadrupole), and we see what is added to the
  record (§2.4).
\end{enumerate}

The same framework is run on four systems (hydrogen atom,
electron--electron, proton--proton, hydrogen molecule), compared with
external measured and literature values, and all programs, data and
figures are stored in reproducible form.

\begin{center}\rule{0.5\linewidth}{0.5pt}\end{center}

\section{2. Framework}\label{framework}

\subsection{2.1 Units and inputs}\label{units-and-inputs}

The units are Gaussian geometric units with \(G=c=q_0=1\), \(q_0=e/3\).
In these units \(\alpha=e^2/(\hbar c)=9q_0^2/(\hbar c)\) fixes
\(\hbar=9/\alpha\), so the Planck constant is not an independent input.
The independent dimensionless inputs are the following four.

{\def\LTcaptype{none} % do not increment counter
\begin{longtable}[]{@{}
  >{\raggedright\arraybackslash}p{(\linewidth - 6\tabcolsep) * \real{0.2500}}
  >{\raggedright\arraybackslash}p{(\linewidth - 6\tabcolsep) * \real{0.2500}}
  >{\raggedright\arraybackslash}p{(\linewidth - 6\tabcolsep) * \real{0.2500}}
  >{\raggedright\arraybackslash}p{(\linewidth - 6\tabcolsep) * \real{0.2500}}@{}}
\toprule\noalign{}
\begin{minipage}[b]{\linewidth}\raggedright
Input
\end{minipage} & \begin{minipage}[b]{\linewidth}\raggedright
Symbol
\end{minipage} & \begin{minipage}[b]{\linewidth}\raggedright
Value
\end{minipage} & \begin{minipage}[b]{\linewidth}\raggedright
Source
\end{minipage} \\
\midrule\noalign{}
\endhead
\bottomrule\noalign{}
\endlastfoot
fine-structure constant & \(\alpha\) & \(7.297352564\times10^{-3}\)
(\(1/137.035999177\)) & CODATA 2022 {[}1{]} (H₂: 2018 {[}2{]},
\(7.2973525693\times10^{-3}\)) \\
mass ratio & \(\mu_p=m_p/m_e\) & 1836.152673426 (H₂: 1836.15267343) &
same \\
strength of gravity &
\(\mu=m_e\sqrt{G}/q_0=3\sqrt{\alpha_{G,e}/\alpha}\) &
\(1.469880\times10^{-21}\)
(\(\alpha_{G,e}=Gm_e^2/\hbar c=1.751809\times10^{-45}\)) & from the SI
values of \(G\), \(m_e\), \(\hbar\), \(c\) {[}1{]} \\
proton \(g\)-factor & \(g_p\) & 5.5856946893 & {[}1{]} \\
\end{longtable}
}

Gravity enters only through the single ratio
\(g=Gm_em_p/e^2=\mu_p\mu^2/9=4.407886\times10^{-40}\). \(g_p\) is the
only external input due to the internal structure of the proton; the
electron has the Dirac value \(g=2\) (no QED). For the hydrogen molecule
the literature Born--Oppenheimer potential \(E_{\rm el}(R)\) {[}46,
47{]}, the adiabatic correction \(E_a(R)\) {[}48{]}, the quadrupole
moment \(\Theta(R)\) {[}49{]} and the rotational \(g\)-factor \(g(R)\)
{[}50{]} are further inputs (§5.1).

Outputs are displayed in eV and seconds, but only ratios are physical.
The energy scale \(m_ec^2=510998.95069\) eV is used only for conversion
in the display.

\subsection{2.2 The rest frame of the absorber and the meaning of ``no
background
spacetime''}\label{the-rest-frame-of-the-absorber-and-the-meaning-of-no-background-spacetime}

The state consists only of the occupations of levels (or of \(m\)
sublevels) and the centre-of-mass velocity \(\vec v\) relative to the
absorber; it holds no coordinate grid and no metric. Emitted photons and
gravitational waves are absorbed by a complete absorber that covers all
directions without gaps at a place far enough from the pair that neither
differences of viewing angle nor distances can be measured. The rest
frame of the absorber is the only reference frame and corresponds to the
asymptotic region in which the momentum and the centre of mass of an
isolated system are defined in general relativity (ADM, Bondi {[}60,
61{]}). The formulas of standard theory (Einstein A coefficients, the
1PN expansion) presuppose an asymptotically flat spacetime; this paper
places that presupposition in this rest frame.

A uniform translation and the position of the centre of mass cannot be
read in principle. What can be read is which cell of the absorber
received the photon (the direction) and its frequency (Doppler). The
number of absorber cells is set to \(N=10^{60}\); directions are treated
as continuous and only the resolution floor
\(\sqrt{4\pi/N}=3.5\times10^{-30}\) rad is kept. It has been checked
that this value does not appear in the record (from the relative width
\(3.2\times10^{-25}\) of the sharpest line, 21 cm, and the recoil
\(v/c=1.1\times10^{-8}\), the upper limit at which the angular
resolution would limit the record is \(N\approx1.5\times10^{34}\)).

\subsection{2.3 The principle of
readout}\label{the-principle-of-readout}

The only information that leaves the system is the events of emission
and absorption. The record of each event is (time, kind of quantum,
frequency, single or pair, direction, polarization). Readouts are made
from this record alone; the level occupations and the dwell times
\(\tau\) are computed as bookkeeping for checking the calculation, but
are not called observables. Among the figures of this paper, the time
evolution of occupations (Fig. 3, Fig. 4 left) is bookkeeping and is
distinguished from emission records (Figs. 2, 5, 10).

\subsection{2.4 The three layers of the
absorber}\label{the-three-layers-of-the-absorber}

The properties given to the absorber are divided into three layers and
introduced in the order 2 → 1 → 3.

\begin{itemize}
\tightlist
\item
  \textbf{Layer 2 (temperature)}: the absorber is a radiation field with
  the Planck distribution \(\bar n(\omega)=1/(e^{\hbar\omega/kT}-1)\);
  to each edge stimulated emission \(A(1+\bar n)\) and absorption
  \(A\bar n\,g_i/g_f\) are added (Einstein 1917 {[}28{]}, \(g=2F+1\) or
  \(2J+1\)). No temperature is applied to gravitational waves or to the
  two-photon channel. At \(T=0\) the isotropic complete absorber is
  recovered.
\item
  \textbf{Layer 1 (direction and velocity)}: a mark \(\hat z\) (a dipole
  axis) is placed in the absorber's rest frame, and the centre of mass
  of the pair is given a finite inertia \(M\) and a velocity \(\vec v\).
  Each emission or absorption changes \(\vec v\) by momentum
  conservation (recoil), and the recorded frequency is the exact Doppler
  \(\omega_{\rm lab}=\omega'/[\gamma(1-\beta n_\parallel)]\). In the
  rest frame of the atom the background becomes dipolar,
  \(T(\hat n)=T_0/[\gamma(1-\beta\hat n\cdot\hat z)]\), and absorption
  and stimulated emission acquire a directional bias. \(v_0=370\) km/s
  matches the velocity of the solar-system barycentre read from the
  dipole of the background radiation, 369.82 ± 0.11 km/s {[}42{]}.
\item
  \textbf{Layer 3 (alignment)}: the basis is resolved down to \(m\)
  sublevels, and the rates are split by Wigner--Eckart as
  \(A_{m\to m'}=A\,(2J+1)\begin{pmatrix}J'&k&J\\ m'&q&-m\end{pmatrix}^2\).
  The background is given a quadrupole \(a_2\),
  \(T(\hat n)=T_0/[\gamma(1-\beta x)]\,(1+a_2P_2(x))\), and the
  sublevels are pumped by the occupation number of each type
  \(\Delta m=q\), \(\langle\bar n\rangle_q=\int\bar n\,P_q\,d\Omega\).
  \(a_2=4\times10^{-6}\) is a guide value, the COBE DMR quadrupole
  \(Q_{\rm rms}=10.7\,\mu\)K {[}43{]} divided by \(T_0\); the actual
  axis of the quadrupole differs from that of the dipole, but here both
  are placed on the same axis.
\end{itemize}

\subsection{2.5 The geometry of the absorber: a map onto the null
surface of the light
cone}\label{the-geometry-of-the-absorber-a-map-onto-the-null-surface-of-the-light-cone}

The absorber of §2.2 is restated in spacetime (Fig. 16). The information
of one emission event lies only on the future light cone of that event.
The absorber at infinity is the cut of this cone by future null
infinity, a different cut for each event. The record (retarded time
\(u\), direction \(\hat n\), frequency, polarization) consists of
coordinates on this cut, and the free data on the whole family of cuts
\(\mathbb R_u\times S^2\) are all the degrees of freedom of the
radiation field (Bondi--Sachs {[}61{]}). In this sense ``the record of
the absorber'' and ``the radiation field'' name the same thing.

Cut by the time slice \(t=t_0\) of the absorber's rest frame the cut is
a sphere; cut by the time slice \(t=t_0+\beta x\) of the rest frame of a
pair moving with velocity \(\beta\) it is the ellipsoid

\[
r(\theta)=\frac{t_0}{1-\beta\cos\theta}.
\]

The background temperature \(T(\hat n)=T_0/[\gamma(1-\beta\cos\theta)]\)
of layer 1 and the aberration (Appendix B) are exactly the geometry of
this tilted cut. The ellipsoid is not the shape of the absorber but the
shape of the same cone cut by a different time slice.

The past and future light cones share only their vertex (the event). The
incoming background radiation (2.7255 K) is the cut of the past cone by
the last-scattering surface; the outgoing record is the cut of the
future cone by infinity (in the universe, the event horizon); they are
different surfaces (§6.5).

Along a null geodesic the phase of a wave is constant (\(k\cdot dx=0\)
gives \(d\phi=0\)). A photon has no proper time, and the phase relation
at emission reaches the cut of the absorber unchanged. The Doppler
factor is not a change of phase along the way but the mismatch between
the clocks of the source and of the absorber at the two ends of the
cone. Therefore the phase relations between records are defined even
though no time evolution of the field is held in the state.

This map is the same operation that has appeared repeatedly in this
research series. The central projection
\(\pi:\mathbb R^n\to S^{n-1}(r_1)\) {[}S4{]} and the radial projection
\(\sigma_R:\mathbb R^{n+1}\setminus\{0\}\to S^n(R)\) {[}S5{]} are
projections from the vertex onto the cut of the cone (the celestial
sphere), and the readout of a wave by amplitude and frequency alone is
likewise the operation of mapping the event at the vertex onto
(direction, frequency, phase) on the cut. The absorber of this paper is
that operation with the time \(u\) of the cut and the frame that cuts it
(the absorber's rest frame) attached. ``No background spacetime'' and
``a complete absorber at infinity'' do not contradict each other; the
latter is a map onto the null surface of the light cone.

\pandocbounded{\includegraphics[keepaspectratio,alt={Fig. 16}]{figs/fig16_lightcone_absorber.png}}
\textbf{Fig. 16.} The absorber is a map of the null surface of the light
cone. Left: the future and past light cones with the emission event at
the vertex. Cut by the time slice of the absorber's rest frame, a sphere
(blue); cut by that of the two-body rest frame, the ellipsoid
\(r(\theta)=t_0/(1-\beta\cos\theta)\) (red). The cut of the past cone
(orange) is the last-scattering surface. Phase marks on the null
generators keep their spacing. Right: the record of each event emitted
from the two-body worldline lands on a separate cut \((u,\hat n)\) of
future infinity (the horizon). Between events there is no readable
change.

\begin{center}\rule{0.5\linewidth}{0.5pt}\end{center}

\section{3. System 1: the hydrogen atom (proton and
electron)}\label{system-1-the-hydrogen-atom-proton-and-electron}

\subsection{3.1 Levels}\label{levels}

The basis is \((n,\ell,s_e,s_p)\), \(n\le5\), with \(s_e\) the sign of
\(j-\ell\) and \(s_p\) the sign of \(F-j\): 50 levels. The level energy
is

\[
E = E_{\rm C} + E_{\rm G} + E_{\rm HFS}.
\]

\textbf{Coulomb (electromagnetic).} To the exact Dirac--Coulomb solution
\(f(n,j)\) of the electron {[}4, 5{]} are added the two-body recoil
\(E_M\) (Barker--Glover {[}6{]}, CODATA 2010 {[}3{]} eq. (16)),

\[
E_M=-\frac{\mu^2(f-1)^2}{2(m_e+m_p)}+\frac{\alpha^4\mu^3}{2n^3m_p^2}\left[\frac{1}{j+\tfrac12}-\frac{1}{\ell+\tfrac12}\right](1-\delta_{\ell0}),
\]

and the higher orders of the retarded photon exchange \(E_S\)
(Salpeter's \((Z\alpha)^5m^2/M\) {[}8, 9, 10{]}, with the Bethe
logarithms \(\ln k_0(n,\ell)\) of Drake--Swainson {[}17{]} and
Jentschura--Mohr {[}18{]}), \((Z\alpha)^6m^2/M\) (\(D_{60}=4\ln2-7/2\)
for S states {[}11, 12{]},
\([3-\ell(\ell+1)/n^2]\cdot2/((4\ell^2-1)(2\ell+3))\) for \(\ell\ge1\)
{[}13, 14{]}), and the squared-logarithm term of order \((Z\alpha)^7\)
(\(D_{72}=-11/(60\pi)\) {[}15, 16{]}), in the form of CODATA 2010 eqs.
(28)--(33). QED (the Lamb shift) and the nuclear size are not included.

\textbf{Gravity.} The two-body 1PN Hamiltonian of Barker--O'Connell
{[}22{]} is quantized in the ordering fixed by the Foldy--Wouthuysen
transformation {[}24{]} of the Dirac Hamiltonian in a static metric,
\(H=\beta mV+\tfrac12\{F,\boldsymbol\alpha\cdot\mathbf p\}\) (Obukhov
{[}23{]}),

\[
m\Phi+\frac{3}{4m}\{p^2,\Phi\}+\frac{3\hbar^2}{8m}\nabla^2\Phi+\frac{3}{2m}(\nabla\Phi\times\mathbf p)\cdot\mathbf S,
\]

and the Newtonian term, the EIH self and cross terms, the \(G^2\) term,
the gravity--electromagnetism 1PN cross term
\(+\tfrac{27}{2}G(m_e+m_p)/r^2\) (Khalil et al.~2018 {[}25{]}; EFT
re-derivation Gupta 2025 {[}26{]}, eq. (4.8)) and the spin--orbit term
\((G/r^3)[(2+\tfrac32m_p/m_e)\mathbf L\cdot\mathbf S_e+(2+\tfrac32m_e/m_p)\mathbf L\cdot\mathbf S_p]\)
are placed on the levels as expectation values. For 1s this gives
\(-1.1988\times10^{-38}\) eV.

\textbf{Hyperfine.}
\(E_{\rm HFS}=K[F(F+1)-j(j+1)-\tfrac34]/[4n^3(\ell+\tfrac12)j(j+1)]\),
\(K=g_p\alpha^4\mu^3/\mu_p\) (Fermi {[}20{]}). The relativistic
corrections (\((3/2)\alpha^2\) for 1s, \((17/8)\alpha^2\) for 2s
{[}7{]}) were confirmed by the numerical solver of §3.6.

\subsection{3.2 Emission rates and time
evolution}\label{emission-rates-and-time-evolution}

Einstein A coefficients are computed for all downward pairs.

\[
A_{E1}=\frac{4\alpha\omega^3}{3c^2}\,|\langle n'\ell'|r|n\ell\rangle|^2\frac{\ell_>}{2\ell+1}\,\mathcal R_1,
\]

\[
A_{E2}=\frac{\alpha\omega^5}{15c^4}\,|\langle r^2\rangle|^2(2\ell'+1)\begin{pmatrix}\ell&2&\ell'\\0&0&0\end{pmatrix}^2\mathcal R_2
\]

\(\mathcal R_k\) is the recoupling to \(j\) and \(F\) (6j symbols,
Racah's formula {[}31{]}, exact in rational arithmetic). M1 is taken for
all pairs within the same \(n\ell\) (hyperfine \(\Delta F=\pm1\) and
fine-structure \(\Delta j=\pm1\)) with the reduced matrix elements of
\(\boldsymbol\mu=-\mu_B(\mathbf L+2\mathbf S)+g_p\mu_N\mathbf I\)
(Edmonds 7.1.7 / 7.1.8 {[}30{]}). For 2s→1s the two-photon rate 8.2206
s⁻¹ (Breit--Teller {[}32{]}, Goldman {[}33{]}) and the relativistic M1
rate \(2.496\times10^{-6}\) s⁻¹ (Johnson {[}34{]}) are given as
constants. Gravitational waves have the matrix elements of the same
quadrupole operator \(r_ir_j\), so the emission ratio to E2 equals the
ratio of the classical quadrupole formulas \((G/45)/(1/180)=4G\)
(Landau--Lifshitz §110 / §71 {[}27{]}):

\[
W_{GW}=\frac{4G\mu^2M^2}{e^2(m_p-m_e)^2}=9.6129\times10^{-43},
\]

\(A_{GW}=W_{GW}A_{E2}\). The numbers of edges are E1 320, E2 376, GW
376, M1 59, two-photon 2.

The time evolution is the master equation
\(\dot{\mathbf p}=R\mathbf p\). The dwell times
\(\boldsymbol\tau=-R_{TT}^{-1}\mathbf p_0\) (transient set \(T\)) are
solved exactly as a linear system, and the emission is counted by
multiplying the flow \(A\tau_{\rm src}\) of each edge by \(\Delta E\)
(no time-step error). Occupations as functions of time are integrated
with the implicit Radau method. The initial state is 5p₃/₂ F=1.

\subsection{3.3 Comparison with external
values}\label{comparison-with-external-values}

{\def\LTcaptype{none} % do not increment counter
\begin{longtable}[]{@{}
  >{\raggedright\arraybackslash}p{(\linewidth - 6\tabcolsep) * \real{0.2500}}
  >{\raggedright\arraybackslash}p{(\linewidth - 6\tabcolsep) * \real{0.2500}}
  >{\raggedright\arraybackslash}p{(\linewidth - 6\tabcolsep) * \real{0.2500}}
  >{\raggedright\arraybackslash}p{(\linewidth - 6\tabcolsep) * \real{0.2500}}@{}}
\toprule\noalign{}
\begin{minipage}[b]{\linewidth}\raggedright
Quantity
\end{minipage} & \begin{minipage}[b]{\linewidth}\raggedright
This work
\end{minipage} & \begin{minipage}[b]{\linewidth}\raggedright
Reference
\end{minipage} & \begin{minipage}[b]{\linewidth}\raggedright
Ratio / difference
\end{minipage} \\
\midrule\noalign{}
\endhead
\bottomrule\noalign{}
\endlastfoot
1S ionization energy & 13.5984684 eV & measured 13.5984346 eV {[}35,
37{]} & difference 8170.31 MHz. Of the 1S Lamb shift 8172.84, the recoil
2.40 is included here, so the remainder should be QED + nuclear size
8170.44. Mismatch 0.12 MHz \\
higher orders of retardation 1S / 2S {[}kHz{]} & \(E_S\) 2409.51 /
341.29, \((Z\alpha)^6\) −7.39 / −0.92, \((Z\alpha)^7\log^2\) −0.42 /
−0.05 & Eides--Grotch--Shelyuto {[}19{]} Table VIII: 2409.51 / 341.29,
−7.38 / −0.92, −0.42 / −0.05 & all digits agree \\
2P₃/₂ − 2P₁/₂ & 10943.68 MHz & measured 10969.04 MHz {[}38{]} &
difference 25.4 MHz \(=(\alpha/\pi)\times\) splitting (QED) \\
2S₁/₂ − 2P₁/₂ & 0.35 MHz (difference of recoil) & measured 1057.85 MHz &
the remaining 1057.5 MHz is QED (Lamb) \\
\(g=Gm_em_p/e^2\) & \(4.407886\times10^{-40}\) & independently from SI
\(4.407886\times10^{-40}\) & 1.00000000 \\
10 E1 A coefficients (2p→1s \(6.2651\times10^8\) s⁻¹ etc.) & & NIST ASD
{[}35, 36{]} & max difference \(1.2\times10^{-4}\) \\
hyperfine 1s / 2s / 2p₁/₂ / 2p₃/₂ & 1418.84 / 177.36 / 59.12 / 23.65 MHz
& measured 1420.41 / 177.56 / 59.22 / 23.65 {[}40{]} & 0.9989
\(=1/(g_e/2)\) (QED) \\
A of 21 cm & \(2.868\times10^{-15}\) s⁻¹ & measured
\(2.884\times10^{-15}\) & 0.9944 (\((g_e/2)^2\) and \(\omega^3\)) \\
E2 / GW of 3d→1s & 593.8 s⁻¹ / \(5.71\times10^{-40}\) s⁻¹ & literature
594 s⁻¹ & \\
identities of expectation values, sum rules, projection theorem & & &
\(10^{-16}\) \\
\end{longtable}
}

Every mismatch equals the size of the QED and nuclear-size terms that
were not included.

\subsection{3.4 Results}\label{results}

\textbf{Emission record} (Figs. 1, 2). The total emission 13.054539268
eV equals the level drop 13.054539268 eV (difference
\(3.6\times10^{-15}\) eV). Breakdown:

{\def\LTcaptype{none} % do not increment counter
\begin{longtable}[]{@{}ll@{}}
\toprule\noalign{}
Kind & Emission {[}eV{]} \\
\midrule\noalign{}
\endhead
\bottomrule\noalign{}
\endlastfoot
E1 & 11.84379 \\
two-photon (2s→1s) & 1.210747 (fraction passing through 2s: 11.9 \%) \\
M1 & \(2.410\times10^{-6}\) \\
E2 & \(1.477\times10^{-6}\) \\
gravitational waves & \(1.420\times10^{-48}\) \\
\end{longtable}
}

The final state is 1s F=0 with probability 1.0 (\(t=10^{17}\) s; 1s F=1
has a 21 cm lifetime of \(1.1\times10^7\) years). The increment of
gravitational binding is \(-1.151\times10^{-38}\) eV \(=2g\times13.05\)
eV.

\textbf{Time structure} (Fig. 3). Within \(10^{-8}\) s 88 \% falls to
1s; 12 \% stays in 2s and decays by two-photon emission in 0.1 s.

\textbf{Widths.} The line width is \(\Gamma_i+\Gamma_f\): Lyman-α 99.71
MHz (measured 99.7), coherence length \(c/(\Gamma_i+\Gamma_f)=0.48\) m.
For the two-photon pair only the sum frequency is sharp
(\(\Gamma_{2s}+\Gamma_{1s}=8.2\) s⁻¹); each photon is continuous
(coherence time \(6.5\times10^{-17}\) s, 19 nm). The spatial extent of
the levels is \(\langle r\rangle=79.4\) pm for 1s, with
\(\Delta r\cdot\Delta p/\hbar=0.866\).

\pandocbounded{\includegraphics[keepaspectratio,alt={Fig. 1}]{figs/fig01_H_levels_cascade.png}}
\textbf{Fig. 1.} Hydrogen atom: 50 levels (hyperfine summed) and the
cascade from 5p₃/₂ F=1. Arrow width is proportional to the square root
of the number of photons per cascade. Solid E1, dashed two-photon,
dotted M1 and E2.

\pandocbounded{\includegraphics[keepaspectratio,alt={Fig. 2}]{figs/fig02_H_spectrum_widths.png}}
\textbf{Fig. 2.} Emission record of the hydrogen atom (\(T=0\)). Top:
expected photon number versus \(\hbar\omega\), coloured by kind. Bottom:
line width \(\Gamma=\Gamma_{\rm src}+\Gamma_{\rm dst}\) and coherence
length \(c/\Gamma\).

\pandocbounded{\includegraphics[keepaspectratio,alt={Fig. 3}]{figs/fig03_H_time_evolution.png}}
\textbf{Fig. 3.} Time evolution of the occupation of each shell \(n\)
(bookkeeping) and the cumulative emission by kind.

\subsection{3.5 Readouts with the three layers of the
absorber}\label{readouts-with-the-three-layers-of-the-absorber}

\textbf{Layer 2 (temperature 2.7255 K)} (Fig. 4).
\(kT=2.3487\times10^{-4}\) eV. For 21 cm,
\(\hbar\omega=5.8679\times10^{-6}\) eV, \(\hbar\omega/kT=0.02498\),
\(\bar n=39.53\). For Lyman-α, \(\hbar\omega/kT=4.3\times10^4\) and
\(\bar n=0\). The 1s F=1→0 rates are spontaneous
\(2.868\times10^{-15}\), stimulated \(1.134\times10^{-13}\), absorption
0→1 \(3.401\times10^{-13}\) s⁻¹, relaxation rate \(4.563\times10^{-13}\)
s⁻¹ (time constant \(2.191\times10^{12}\) s \(=6.94\times10^4\) years).
The steady distribution is F=1 : F=0
\(=0.745286:0.254714=3e^{-\hbar\omega/kT}\) (difference
\(4\times10^{-16}\)), steady residual \(\max|R\mathbf p|=0\). The spin
temperature read from the ratio of the 21 cm doublet is 2.72550 K, equal
to the absorber temperature. The multistep emission from 5p is
unchanged.

\pandocbounded{\includegraphics[keepaspectratio,alt={Fig. 4}]{figs/fig04_H_layer2_temperature.png}}
\textbf{Fig. 4.} Layer 2. Left: time evolution of the occupation of the
1s hyperfine doublet (\(T=0\) and 2.7255 K). Right: temperature
dependence of the ratio \(p(F{=}1)/p(F{=}0)\) at \(t=10^{17}\) s and
\(3e^{-\hbar\omega/kT}\).

\textbf{Layer 1 (direction, centre-of-mass velocity, recoil)} (Fig. 5).
The recoil shift of Lyman-α \(\Delta E^2/(2Mc^2)=13.40\) MHz, the recoil
velocity per photon \(\hbar\omega/(Mc)=3.257\) m/s, for 21 cm
\(1.874\times10^{-6}\) m/s. First-order Doppler
\(\beta=1.2342\times10^{-3}\), second order
\(\beta^2/2=7.616\times10^{-7}\). The background temperature seen in the
atom's rest frame is 2.728866 K forward and 2.722138 K backward.
Readouts:

\begin{enumerate}
\def\labelenumi{\arabic{enumi}.}
\tightlist
\item
  From the record \((\delta,\hat n)\) of 212 multistep photons of an
  ensemble of 200 atoms, a least-squares fit with the exact Doppler
  model gives \(\vec v=(-0.03,-0.03,370000.08)\) m/s, a difference of
  0.09 m/s from the given \((0,0,370000)\) (the same computation in the
  engine version with another random sequence gives 0.11 m/s). The
  first-order linear solution is off by tens of m/s because of the
  second-order Doppler term.
\item
  For the sequence of 21 cm events of a single atom (absorption,
  stimulated and spontaneous emission), the \(\delta\) of each event
  agrees to double precision with the exact Doppler of the true velocity
  (the relative width \(5\times10^{-23}\) of 21 cm is below the
  double-precision floor \(10^{-19}\)).
\item
  From the same 29 events, accumulating the kicks
  \(\mp\hbar\omega\hat n/c\) contained in the record (whether absorption
  or emission is in the record) and solving jointly, the three
  components of \(\vec v_0\) before the first event are recovered to
  \(2.39\times10^{-10}\) m/s (\(2\times10^{-4}\) of the recoil
  \(1.9\times10^{-6}\) m/s).
\end{enumerate}

\pandocbounded{\includegraphics[keepaspectratio,alt={Fig. 5}]{figs/fig05_H_layer1_doppler_recoil.png}}
\textbf{Fig. 5.} Layer 1. Left: \((\delta/\beta,\ \hat n\cdot\hat v_0)\)
of the photons of 200 atoms and the exact Doppler. Middle: residuals
(converted to m/s). Right: the history of \(|\vec v|\) over the 21 cm
events of one atom (true values) and its reconstruction from the record.

\textbf{Layer 3 (alignment, polarization)} (Fig. 6). Basis of 220
states, 15103 \(m\)-resolved edges, maximum relative error of the sum
rule \(\sum_{m'}A_{mm'}=A\) of \(5.5\times10^{-16}\). (A) From the
initial 5p₃/₂ F=1 with \(m_F=0\), the 5p→1s F0 line has \(\Delta m=0\)
only (pure π, \(I\propto\sin^2\theta\), linear polarization degree at
90° equal to +1); from \(m_F=+1\) it is σ; averaged over \(m\) all lines
return to isotropic and unpolarized, and the total emission is
13.054539268 eV. The initial \(m\) appears in the record as the angular
distribution and polarization of the emitted light. (B) Steady alignment
\(A_2=(p_++p_--2p_0)/\sum p\) at 2.7255 K:

{\def\LTcaptype{none} % do not increment counter
\begin{longtable}[]{@{}
  >{\raggedright\arraybackslash}p{(\linewidth - 8\tabcolsep) * \real{0.2000}}
  >{\raggedright\arraybackslash}p{(\linewidth - 8\tabcolsep) * \real{0.2000}}
  >{\raggedright\arraybackslash}p{(\linewidth - 8\tabcolsep) * \real{0.2000}}
  >{\raggedright\arraybackslash}p{(\linewidth - 8\tabcolsep) * \real{0.2000}}
  >{\raggedright\arraybackslash}p{(\linewidth - 8\tabcolsep) * \real{0.2000}}@{}}
\toprule\noalign{}
\begin{minipage}[b]{\linewidth}\raggedright
Background
\end{minipage} & \begin{minipage}[b]{\linewidth}\raggedright
\(\langle\bar n\rangle_0\) (π)
\end{minipage} & \begin{minipage}[b]{\linewidth}\raggedright
\(\langle\bar n\rangle_{\pm1}\) (σ)
\end{minipage} & \begin{minipage}[b]{\linewidth}\raggedright
alignment \(A_2\) of 1s F=1
\end{minipage} & \begin{minipage}[b]{\linewidth}\raggedright
polarization degree of 21 cm emission at 90°
\end{minipage} \\
\midrule\noalign{}
\endhead
\bottomrule\noalign{}
\endlastfoot
isotropic & 39.52787894 & 39.52787894 & 0 & 0 \\
dipole β only (quadrupole at second order) & 39.52786065 & 39.52787285 &
\(+5.074\times10^{-9}\) & \(-3.806\times10^{-9}\) \\
dipole β + quadrupole \(a_2=4\times10^{-6}\) & 39.52782863 & 39.52788886
& \(+2.506\times10^{-8}\) & \(-1.880\times10^{-8}\) \\
\end{longtable}
}

An unpolarized intensity dipole produces no alignment at first order;
the quadrupole component \(\beta^2\approx1.5\times10^{-6}\) at second
order and the quadrupole \(a_2\) of the background do. Since the steady
sublevel occupations are
\(p_m\propto\langle\bar n\rangle_q/(1+\langle\bar n\rangle_q)\), the
relative difference \(10^{-6}\) of the occupation numbers is suppressed
by \(1/(\bar n(1+\bar n))\approx1/1600\), and the alignment is of order
\(10^{-8}\). Reading the linear polarization degree \(2\times10^{-8}\)
of 21 cm would require \(10^{16}\) photons.

\pandocbounded{\includegraphics[keepaspectratio,alt={Fig. 6}]{figs/fig06_H_layer3_alignment.png}}
\textbf{Fig. 6.} Layer 3. Left: angular patterns of the \(\Delta m\)
types (the 5p→1s F0 line from \(m_F=0\) is pure π). Right: steady
alignment of 1s F=1 at 2.7255 K versus the background quadrupole
\(a_2\).

\subsection{3.6 Comparison with a numerical solver that takes the sign
of the charge as
input}\label{comparison-with-a-numerical-solver-that-takes-the-sign-of-the-charge-as-input}

A general solver was built separately that produces the levels not from
closed formulas but from a two-sided shooting of the radial Dirac
equation (solved in \(\eta=1-\varepsilon\) to avoid cancellation) and
from expectation values over Schrödinger solutions, and it was compared
with the fixed set under the same conditions (\(n\le5\), 5p₃/₂ F=1,
\(t_{\max}=10^{17}\) s). Over the 50 levels the maximum difference of
the Coulomb levels is \(3.0\times10^{-11}\) eV (relative \(10^{-11}\)),
the maximum relative difference of gravity \(3.9\times10^{-9}\), the
hyperfine Dirac/NR − 1 is \((3/2)\alpha^2=8.0\times10^{-5}\) for 1s and
2s (reproducing Breit's relativistic correction), and the maximum
difference of the total energy \(3.3\times10^{-10}\) eV. The A
coefficients agree in the set of edges (E1 320, E2 376, M1 58) with a
median relative difference of \(6.6\times10^{-5}\) (the \(\alpha^2\)
difference between the Dirac \((G_aG_b+F_aF_b)\) and the Schrödinger
radial integrals), and the total emission 13.054539268 eV agrees. Since
this solver takes the sign of the charge as input, it is used for the
same-sign pairs of §4.

\begin{center}\rule{0.5\linewidth}{0.5pt}\end{center}

\section{4. Systems 2 and 3: electron--electron,
proton--proton}\label{systems-2-and-3-electronelectron-protonproton}

\subsection{4.1 Absence of bound levels}\label{absence-of-bound-levels}

With everything but the particles kept as for hydrogen (basis \(n\le5\),
the level corresponding to the initial 5p₃/₂ F=1, the absorber), the
solver of §3.6 was given same-sign charges. The net coupling is
\(+7.297\times10^{-3}\) (Coulomb) against gravity
\(+1.752\times10^{-45}\) for electron--electron, and gravity/Coulomb
\(=8.09\times10^{-37}\) for proton--proton. The net potential is
positive (repulsive) at every \(r\), so there are exactly zero bound
levels; the model of levels and rates has nothing to represent, emission
0, no final state. This is the set's own answer to the question whether
``levels and radiation alone'' can solve the problem.

\subsection{4.2 Scattering: the Mott cross
section}\label{scattering-the-mott-cross-section}

What Einstein--Maxwell theory predicts for same-sign pairs is
scattering. The intermediate orbit is not observed; what remains in the
record is the distribution of deflection angles, and for identical
particles assuming a real orbit drops the Mott interference term and
gives the wrong answer (the S-matrix standpoint {[}63{]}). The relative
kinetic energy is set to the initial hydrogen level,
\(E=|E(5p_{3/2}F1)|=0.543934\) eV, and the impact parameter for the
spectrum is \(b=L/(\mu v)\) with the \(\ell=1\) angular momentum
\(L=\hbar\sqrt2\).

{\def\LTcaptype{none} % do not increment counter
\begin{longtable}[]{@{}
  >{\raggedright\arraybackslash}p{(\linewidth - 4\tabcolsep) * \real{0.3333}}
  >{\raggedright\arraybackslash}p{(\linewidth - 4\tabcolsep) * \real{0.3333}}
  >{\raggedright\arraybackslash}p{(\linewidth - 4\tabcolsep) * \real{0.3333}}@{}}
\toprule\noalign{}
\begin{minipage}[b]{\linewidth}\raggedright
Quantity
\end{minipage} & \begin{minipage}[b]{\linewidth}\raggedright
electron--electron
\end{minipage} & \begin{minipage}[b]{\linewidth}\raggedright
proton--proton
\end{minipage} \\
\midrule\noalign{}
\endhead
\bottomrule\noalign{}
\endlastfoot
reduced mass \(\mu\) & 0.5000 \(m_e\) & 918.0763 \(m_e\) \\
\(v/c\) & \(2.0634\times10^{-3}\) & \(4.8155\times10^{-5}\) \\
Sommerfeld parameter \(\eta=(q_aq_b/\hbar c)/(v/c)\) & 3.5365 (quantum
regime) & 151.5398 (semiclassical) \\
wave number \(k\), de Broglie wavelength & 0.1414 \(/a_0\), 2.352 nm &
6.0583 \(/a_0\), 0.055 nm \\
distance parameter \(a=q_aq_b/(\mu v^2)\), head-on closest approach
\(2a\) & 25.014 \(a_0\), 2.647 nm & 25.014 \(a_0\), 2.647 nm \\
impact parameter \(b\), eccentricity \(e\) & 10.003 \(a_0\), 1.0770 &
0.233 \(a_0\), 1.0000 \\
Darwin (magnetic) term \(v^2/c^2\) & \(4.26\times10^{-6}\) &
\(2.32\times10^{-9}\) \\
gravity/Coulomb, gravitational-wave/electric-quadrupole \(4Gm^2/q^2\) &
\(2.401\times10^{-43}\), \(9.602\times10^{-43}\) &
\(8.094\times10^{-37}\), \(3.237\times10^{-36}\) \\
Mott \(d\sigma/d\Omega\) (90°), distinguishable particles both counted &
1.7521, 3.5041 nm²/sr & 1.7521, 3.5041 nm²/sr \\
\end{longtable}
}

The Mott cross section {[}44{]}, obtained by symmetrizing the exact
Coulomb amplitude

\[
f(\theta)=-\frac{\eta}{2k\sin^2(\theta/2)}\exp\!\left[-i\eta\ln\sin^2(\theta/2)+2i\sigma_0\right]
\]

for identical unpolarized spin-½ particles,

\[
\frac{d\sigma}{d\Omega}=\left(\frac{\eta}{2k}\right)^2\left[\frac{1}{\sin^4(\theta/2)}+\frac{1}{\cos^4(\theta/2)}-\frac{\cos(\eta\ln\tan^2(\theta/2))}{\sin^2(\theta/2)\cos^2(\theta/2)}\right],
\]

is one half of the distinguishable-particle value at 90° (fermionic
antisymmetrization). The interference term oscillates more finely the
larger \(\eta\) is (Fig. 7). It diverges in the forward direction and
the total cross section is not defined.

\pandocbounded{\includegraphics[keepaspectratio,alt={Fig. 7}]{figs/fig07_ee_pp_mott.png}}
\textbf{Fig. 7.} Mott cross sections of electron--electron
(\(\eta=3.54\)) and proton--proton (\(\eta=151.5\)) and the classical
distinguishable-particle values.

\subsection{4.3 Quadrupole
bremsstrahlung}\label{quadrupole-bremsstrahlung}

For identical particles the electric dipole in the centre-of-mass frame
vanishes identically, so bremsstrahlung is quadrupole,
\(P=(1/180c^5)\dddot D_{\alpha\beta}\dddot D_{\alpha\beta}\)
(Landau--Lifshitz §71 {[}27{]}). On the repulsive hyperbola
\(x=a(e+\cosh\xi)\), \(y=a\sqrt{e^2-1}\sinh\xi\),
\(t=(a/v)(e\sinh\xi+\xi)\), \(\dddot D\) is obtained in closed form and
Fourier-transformed in \(\xi\) to give \(dE/d\omega\). The Parseval
self-check \(\int P\,dt=\int(dE/d\omega)d\omega\) gives 0.9996 for
electron--electron and 1.0001 for proton--proton.

{\def\LTcaptype{none} % do not increment counter
\begin{longtable}[]{@{}
  >{\raggedright\arraybackslash}p{(\linewidth - 4\tabcolsep) * \real{0.3333}}
  >{\raggedright\arraybackslash}p{(\linewidth - 4\tabcolsep) * \real{0.3333}}
  >{\raggedright\arraybackslash}p{(\linewidth - 4\tabcolsep) * \real{0.3333}}@{}}
\toprule\noalign{}
\begin{minipage}[b]{\linewidth}\raggedright
Quantity
\end{minipage} & \begin{minipage}[b]{\linewidth}\raggedright
electron--electron
\end{minipage} & \begin{minipage}[b]{\linewidth}\raggedright
proton--proton
\end{minipage} \\
\midrule\noalign{}
\endhead
\bottomrule\noalign{}
\endlastfoot
radiated energy \(E_{\rm rad}=\int(dE/d\omega)d\omega\) &
\(1.54\times10^{-16}\) eV & \(3.58\times10^{-25}\) eV \\
characteristic photon energy \(\hbar v/a\) & 0.31 eV &
\(7.18\times10^{-3}\) eV \\
expected photons per collision \(E_{\rm rad}/\hbar(v/a)\) &
\(5.0\times10^{-16}\) & \(5.0\times10^{-23}\) \\
gravitational waves & \(9.6\times10^{-43}\) of that &
\(3.2\times10^{-36}\) of that \\
\end{longtable}
}

The spectrum is flat as \(\omega\to0\) (the difference of \(\ddot D\)
before and after the encounter) and falls exponentially for
\(\omega\gtrsim v/a\) (Fig. 8). For proton--proton, \(\eta\gg1\), the
classical spectrum is valid; for electron--electron, \(\eta\approx3.5\),
it is a guide (no quantum Gaunt factor is included). Nothing comes out
of a single collision; only the in--out relation remains in the record.

\pandocbounded{\includegraphics[keepaspectratio,alt={Fig. 8}]{figs/fig08_ee_pp_bremsstrahlung.png}}
\textbf{Fig. 8.} Spectra \(dE/d\omega\) of quadrupole bremsstrahlung.
The dashed line marks \(\hbar v/a\).

\begin{center}\rule{0.5\linewidth}{0.5pt}\end{center}

\section{5. System 4: the hydrogen molecule (hydrogen atom--hydrogen
atom)}\label{system-4-the-hydrogen-molecule-hydrogen-atomhydrogen-atom}

\subsection{5.1 Potential and levels}\label{potential-and-levels}

Neutron--neutron is meaningless unless the strong force is included; the
two-body systems without charge and with mass that can be built without
assuming the strong force are pairs of neutral atoms, the smallest of
which is H₂. The rovibrational levels \((v,J)\) of the ground electronic
state X¹Σg⁺ are taken as the levels of the system.

The radial nuclear equation in the adiabatic approximation is

\[
-\frac{1}{2\mu_n}f''+\left[V(R)+\frac{J(J+1)}{2\mu_nR^2}\right]f=Ef,\qquad \mu_n=\frac{m_p}{2},\qquad V(R)=E_{\rm el}(R)+E_a(R).
\]

\(E_{\rm el}\) is the analytic fit (Łach; FUNCTION V of H2SPECTRE 7.4
{[}47{]}, ported) of Pachucki's BO potential {[}46{]} (82 points at R =
0.1--20 bohr, precision \(10^{-15}\); the maximum difference between fit
and table is \(2.2\times10^{-12}\) hartree), and \(E_a\) is the fit
(Eaint of the same program) of the adiabatic correction of
Pachucki--Komasa {[}48{]}. Both take two hydrogen atoms as the origin,
\(V(\infty)=0\). Nonadiabatic, relativistic and QED corrections are not
included. The discretization is a sinc-DVR (Colbert--Miller {[}55{]},
the same matrix elements as H2SPECTRE), \(dr=0.025\) bohr, \(R\le50\)
bohr. Changing the grid to \(dr=0.02\), \(R\le60\) changes the levels by
\(2.2\times10^{-7}\) cm⁻¹.

The bound levels are the 301 with \(E<0\) (of the 302 of Roueff et
al.~{[}51{]}, (14,4) is not bound without the nonadiabatic correction;
with all corrections it lies 0.027 cm⁻¹ below dissociation).
\(D_0=36118.3637\) cm⁻¹ (E2(FULL) 36118.7977, all corrections
36118.0695), (0,1)−(0,0) = 118.4919 cm⁻¹ (all corrections 118.4868).
Even \(J\) is para (\(I=0\)), odd \(J\) ortho (\(I=1\)).

\subsection{5.2 Rates}\label{rates}

In atomic units (\(\hbar=m_e=e=1\), \(c=1/\alpha\)),

\[
A_{E2}=\frac{\alpha^5\omega^5}{15}(2J_f+1)\begin{pmatrix}J_i&2&J_f\\0&0&0\end{pmatrix}^2|\langle f_f|\Theta|f_i\rangle|^2,
\]

\[
A_{M1}=\frac{\alpha^5\omega^3}{3}\left(\frac{m_e}{m_p}\right)^2J(J+1)|\langle f_f|g|f_i\rangle|^2.
\]

E2 has \(\Delta J=0,\pm2\); M1 has \(\Delta J=0\), \(v_i\ne v_f\).
\(\Theta(R)\) is the molecular quadrupole moment, one half of
\(Q_{\rm WSD}(R)\) (257 points) of Table 3 of
Wolniewicz--Simbotin--Dalgarno {[}49{]} (the note in §3.1 of Roueff et
al.). In the Q branch
\((2J+1)\begin{pmatrix}J&2&J\\0&0&0\end{pmatrix}^2=J(J+1)/((2J-1)(2J+3))\),
which agrees with eq. (14) of Pachucki--Komasa 2011 {[}50{]}. \(g(R)\)
is Table I of the same paper {[}50{]}. Gravitational waves follow the
same derivation of the ratio as for hydrogen,

\[
A_{GW}=4\frac{\alpha_{Gp}}{\alpha}\frac{|\langle q_m\rangle|^2}{|\langle\Theta\rangle|^2}A_{E2},\qquad q_m(R)=\frac{R^2}{2}+\frac{m_e}{m_p}\left(\frac{R^2}{2}-\Theta(R)\right),\qquad \alpha_{Gp}=\frac{Gm_p^2}{\hbar c}.
\]

The edges are E2 4669, M1 1467, GW 4669. E2, M1 and GW all preserve the
parity of \(J\), so ortho and para are not connected (0 edges change the
parity). The mass is \(M=2(m_p+m_e)+E/c^2\). The nuclear-spin degeneracy
is common to src and dst, so the detailed-balance factor
\(g_{\rm src}/g_{\rm dst}=(2J_s+1)/(2J_d+1)\) can be used as it is.

\subsection{5.3 Validation}\label{validation}

Comparison with Roueff et al.~2019 {[}51{]} (CDS J/A+A/630/A58, 4712
transitions; the same BO + adiabatic \(V\), \(\Theta\) and \(g\) from
Pachucki--Komasa, levels with all corrections from Pachucki--Komasa 2018
{[}52{]}) (Fig. 11).

{\def\LTcaptype{none} % do not increment counter
\begin{longtable}[]{@{}
  >{\raggedright\arraybackslash}p{(\linewidth - 2\tabcolsep) * \real{0.5000}}
  >{\raggedright\arraybackslash}p{(\linewidth - 2\tabcolsep) * \real{0.5000}}@{}}
\toprule\noalign{}
\begin{minipage}[b]{\linewidth}\raggedright
Quantity
\end{minipage} & \begin{minipage}[b]{\linewidth}\raggedright
Result
\end{minipage} \\
\midrule\noalign{}
\endhead
\bottomrule\noalign{}
\endlastfoot
level energies, this work − CDS & −0.29 ((0,0)) to +4.6 cm⁻¹ (\(v=10\));
grows with \(v\) = the nonadiabatic correction \\
relative difference of wavenumbers σ & median \(1.6\times10^{-4}\) (up
to \(4\times10^{-2}\) for small σ between close levels) \\
relative difference of \(A_{E2}\) (4669 lines) & median
\(8.2\times10^{-4}\); \(7.7\times10^{-4}\) with σ⁵ rescaled so that only
the Θ matrix element differs; for strong lines \(A>10^{-9}\) s⁻¹ up to 6
\% (lines with \(\Delta v\ge3\) and large cancellation) \\
relative difference of \(A_{M1}\) (1467 lines) & median
\(5.8\times10^{-4}\) \\
1-0 S(1) (2.12 μm) & \(3.471\times10^{-7}\) vs \(3.470\times10^{-7}\)
s⁻¹ \\
0-0 S(0) (28.2 μm) & \(2.942\times10^{-11}\) vs \(2.943\times10^{-11}\)
s⁻¹ \\
M1 of 1-0 Q(1) & \(7.137\times10^{-10}\) vs \(7.137\times10^{-10}\)
s⁻¹ \\
\end{longtable}
}

The entry (2,24) of the level table H2.dat distributed with H2SPECTRE is
115 cm⁻¹ off from its neighbours in \(J\) and is taken to be a misprint
in the distributed data.

\pandocbounded{\includegraphics[keepaspectratio,alt={Fig. 11}]{figs/fig11_H2_validation.png}}
\textbf{Fig. 11.} Validation of H₂. Left: \(A_{E2}\) of this work
against CDS (4669 lines, coloured by \(\Delta v\)). Middle: distribution
of the relative differences. Right: level-energy differences versus
\(v\).

\subsection{5.4 Results of the cascades}\label{results-of-the-cascades}

With initial levels (1,3) (the upper level of 1-0 S(1)), (3,5), and
(14,0), (14,1) just below the dissociation limit, the same master
equation as in §3.2 was run (Figs. 9, 10).

{\def\LTcaptype{none} % do not increment counter
\begin{longtable}[]{@{}
  >{\raggedright\arraybackslash}p{(\linewidth - 8\tabcolsep) * \real{0.2000}}
  >{\raggedright\arraybackslash}p{(\linewidth - 8\tabcolsep) * \real{0.2000}}
  >{\raggedright\arraybackslash}p{(\linewidth - 8\tabcolsep) * \real{0.2000}}
  >{\raggedright\arraybackslash}p{(\linewidth - 8\tabcolsep) * \real{0.2000}}
  >{\raggedright\arraybackslash}p{(\linewidth - 8\tabcolsep) * \real{0.2000}}@{}}
\toprule\noalign{}
\begin{minipage}[b]{\linewidth}\raggedright
Initial
\end{minipage} & \begin{minipage}[b]{\linewidth}\raggedright
depth below dissociation {[}cm⁻¹{]}
\end{minipage} & \begin{minipage}[b]{\linewidth}\raggedright
total emission {[}eV{]} (\(=E_{\rm init}-E_{\rm final}\), difference)
\end{minipage} & \begin{minipage}[b]{\linewidth}\raggedright
photons E2 / M1
\end{minipage} & \begin{minipage}[b]{\linewidth}\raggedright
final state
\end{minipage} \\
\midrule\noalign{}
\endhead
\bottomrule\noalign{}
\endlastfoot
(1,3) & 31286.1 & 0.584431745 (\(1\times10^{-15}\)) & 1.830 / 0.0051 &
(0,1), ortho ground \\
(3,5) & 22850.8 & 1.630280932 (\(2\times10^{-15}\)) & 3.537 / 0.027 &
(0,1) \\
(14,0) & 143.3 & 4.460335799 (\(2\times10^{-15}\)) & 7.479 / 0.022 &
(0,0), para ground \\
(14,1) & 126.2 & 4.447765200 (\(1\times10^{-14}\)) & 6.972 / 0.024 &
(0,1) \\
\end{longtable}
}

The first step from (14,J) goes to (10,J±2), (9,J±2), (8,J±2)
(\(\hbar\omega\) 0.6--1.1 eV, branching 0.37 / 0.36 / 0.14) and not
directly to \(v'=0\). Thereafter \(v\) decreases by 1--3 at a time, and
the cascade ends with the pure-rotational (0,2)→(0,0) (lifetime
\(3.4\times10^{10}\) s, 28 μm) or (0,3)→(0,1) (lifetime
\(2.1\times10^9\) s, 17 μm, \(A=4.76\times10^{-10}\) s⁻¹). The natural
widths \(\Gamma\hbar/\hbar\omega\) of the lines are
\(10^{-24}\)--\(10^{-20}\), far sharper than the \(10^{-8}\) of the
hydrogen atom. The gravitational-wave emission is
\(1.5\times10^{-35}\)--\(5.4\times10^{-35}\) eV.

\pandocbounded{\includegraphics[keepaspectratio,alt={Fig. 9}]{figs/fig09_H2_levels_cascades.png}}
\textbf{Fig. 9.} The 301 levels of H₂ X¹Σg⁺ (abscissa \(J\), ordinate
the excitation above (0,0)) and the E2 cascades from (1,3) and (14,1).

\pandocbounded{\includegraphics[keepaspectratio,alt={Fig. 10}]{figs/fig10_H2_cascade_spectra.png}}
\textbf{Fig. 10.} Emission records of the cascades (photon number versus
wavelength). From top: (1,3), (3,5), (14,1).

\subsection{5.5 Readouts with the three layers of the
absorber}\label{readouts-with-the-three-layers-of-the-absorber-1}

\textbf{Layer 2.} The lowest excitation from the ground level,
(0,1)→(0,3), is 587 cm⁻¹ \(=0.0728\) eV \(=310\,kT\) (2.7255 K), with
Boltzmann ratio \(g_f/g_i\,e^{-\hbar\omega/kT}=6\times10^{-135}\). The
largest Planck occupation among all edges is
\(\langle\bar n\rangle=0.65\) for the accidentally close pair of highly
excited levels (6,20)→(7,18) (1.77 cm⁻¹). The steady state is the single
lowest level of the nuclear-spin series, and there is no temperature
readout like that of the hydrogen 21 cm line (\(\bar n=39.5\)).

\textbf{Layer 1} (Fig. 13). For the cascade from (1,3), 200 molecules
and 360 photons, the recoil is 0.01--0.09 m/s and the line widths
\(2\times10^{-24}\)--\(7\times10^{-22}\). Recovering \(\vec v\) from the
record \((\delta,\hat n)\) alone gives a difference of 0.002 m/s (the
first-order linear solution is off by 27 m/s). For (3,5), 695 photons
give 0.007 m/s; for (14,1), 1430 photons give 0.026 m/s (recoil up to
0.32 m/s). No absorption or stimulated emission occurs within
\(t\le10^{17}\) s.

\textbf{Layer 3} (Fig. 14). From the initial \(m=0\), the 1-0 S(1) line
((1,3)→(0,1)) has \(N_0:N_{+1}:N_{-1}=0.6:0.2:0.2\) (polarization degree
at 90° equal to +0.5), and the M1 line 1-0 Q(3) has zero \(\Delta m=0\)
component (\((3\,1\,3;000)=0\)). The \(m\) distribution of the final
state (0,1) is 0.240 / 0.521 / 0.240 (alignment −0.187); starting from
(14,1) it is 0.330 / 0.339 / 0.330 (−0.006). Since the background has
\(\langle\bar n\rangle\approx0\), no relaxation or alignment between
\(m\) sublevels occurs, and the memory of the initial \(m\) remains in
the final state as it is.

\pandocbounded{\includegraphics[keepaspectratio,alt={Fig. 13}]{figs/fig13_H2_layer1_doppler.png}}
\textbf{Fig. 13.} Layer 1 of H₂. Doppler record of the 360 photons from
(1,3) and the residuals of the recovery of \(\vec v\) with the exact
model.

\pandocbounded{\includegraphics[keepaspectratio,alt={Fig. 14}]{figs/fig14_H2_layer3_polarization.png}}
\textbf{Fig. 14.} Layer 3 of H₂. Distribution of \(\Delta m\) of each
line from the initial \(m=0\), and the \(m\) distribution of the final
state (0,1).

\subsection{5.6 Radiative association H + H → H₂ +
γ}\label{radiative-association-h-h-hux2082-ux3b3}

Two hydrogen atoms do not approach each other by gravity. The
inter-atomic gravitational potential \(-Gm_H^2/R\) is
\(8\times10^{-37}\) of the electromagnetic interaction; it exceeds the
dispersion force (Casimir--Polder \(-C_7/R^7\)) only for \(R\gtrsim0.2\)
mm, but to be gravitationally bound there would require a
``gravitational atom'' with Bohr radius \(7\times10^{22}\) m (2.3 Mpc)
and binding energy \(8\times10^{-69}\) eV. The actual path splits by
electron spin: the triplet b³Σu⁺ (3/4) is repulsive, while the singlet
X¹Σg⁺ (1/4) falls into a well 4.75 eV deep, but the nuclei approach only
to the inner turning point (\(R=0.776\,a_0\) at the dissociation limit)
under the Coulomb repulsion and return. To unite, 4.5 eV must be given
to a partner; with two bodies only, radiation is the only way, and E1 is
absent for a homonuclear pair.

Therefore the rates from discretized continuum states \(\psi_n\)
(\(E_n>0\), partial wave \(J\)) normalized in a box \(R\le50\) bohr to
the bound levels were computed with the same formulas as in §5.2, and

\[
P_J(E_n)=2\pi\hbar\rho_J(E_n)\sum_fA(n\to f),\qquad
\sigma(E)=\frac{\pi}{k^2}\sum_Jw_J(2J+1)P_J(E),\qquad k(E)=\sigma v,
\]

where \(\rho_J=dn/dE\) is the level density of the box and
\(2\pi\hbar\rho_J\) is the round-trip time in the box
(\(2.181\times10^{-12}\) s for \(J=0\), \(E\approx100\) cm⁻¹; the
classical \(2\int dR/v\) is \(2.161\times10^{-12}\) s). \(P_J\) is the
capture probability per pass and is independent of \(R_{\max}\)
(relative difference \(10^{-5}\) between 50 and 100 bohr). The weights
\(w_J\) are the electronic singlet 1/4 times the nuclear-spin weights
(1/4 for even \(J\), 3/4 for odd \(J\)). \(J\le24\), \(E\le4389\) cm⁻¹.

{\def\LTcaptype{none} % do not increment counter
\begin{longtable}[]{@{}
  >{\raggedright\arraybackslash}p{(\linewidth - 8\tabcolsep) * \real{0.2000}}
  >{\raggedright\arraybackslash}p{(\linewidth - 8\tabcolsep) * \real{0.2000}}
  >{\raggedright\arraybackslash}p{(\linewidth - 8\tabcolsep) * \real{0.2000}}
  >{\raggedright\arraybackslash}p{(\linewidth - 8\tabcolsep) * \real{0.2000}}
  >{\raggedright\arraybackslash}p{(\linewidth - 8\tabcolsep) * \real{0.2000}}@{}}
\toprule\noalign{}
\begin{minipage}[b]{\linewidth}\raggedright
\(E\) {[}cm⁻¹{]}
\end{minipage} & \begin{minipage}[b]{\linewidth}\raggedright
\(E/k_B\) {[}K{]}
\end{minipage} & \begin{minipage}[b]{\linewidth}\raggedright
\(P_0\)
\end{minipage} & \begin{minipage}[b]{\linewidth}\raggedright
\(\sigma\) {[}cm²{]}
\end{minipage} & \begin{minipage}[b]{\linewidth}\raggedright
\(k(E)\) {[}cm³ s⁻¹{]}
\end{minipage} \\
\midrule\noalign{}
\endhead
\bottomrule\noalign{}
\endlastfoot
3 & 4.3 & \(9.4\times10^{-20}\) & \(1.12\times10^{-34}\) &
\(4.2\times10^{-30}\) \\
10 & 14 & \(1.03\times10^{-19}\) & \(8.9\times10^{-35}\) &
\(6.1\times10^{-30}\) \\
100 & 144 & \(1.12\times10^{-19}\) & \(9.8\times10^{-35}\) &
\(2.1\times10^{-29}\) \\
1000 & 1439 & \(1.18\times10^{-19}\) & \(4.4\times10^{-35}\) &
\(3.1\times10^{-29}\) \\
3000 & 4316 & \(1.27\times10^{-19}\) & \(2.3\times10^{-35}\) &
\(2.7\times10^{-29}\) \\
\end{longtable}
}

The Maxwell--Boltzmann average \(k(T)\) is \(2.4\times10^{-29}\) at 10
K, \(3.2\times10^{-29}\) at 100 K, \(3.8\times10^{-29}\) at 300 K,
\(3.5\times10^{-29}\) at 1000 K and \(2.4\times10^{-29}\) cm³ s⁻¹ at
3000 K. The main final levels are \(v'=8\)--11 (\(\hbar\omega\) 0.4--1.1
eV); capture does not go to \(v'=0\). At \(E\approx1\) cm⁻¹ the
quasi-bound level (14,4) (0.74 cm⁻¹ above the dissociation limit in this
calculation, \(P_4=9\times10^{-18}\)) accounts for 99 \% of σ; with the
nonadiabatic correction it returns to a bound level and disappears from
the continuum. The sharp peaks in \(k(E)\) are quasi-bound \((v,J)\)
behind the centrifugal barrier (shape resonances), whose widths the box
discretization does not resolve.

The estimate made before the calculation, \(10^{-14}\) from the E2 rate
at \(\hbar\omega=D_e\) and the transit time through the well, was five
orders of magnitude too large. The actual capture is a transition with
\(\hbar\omega\approx0.4\)--1.1 eV, for which \(\omega^5\) is \(10^{-4}\)
times smaller.

\pandocbounded{\includegraphics[keepaspectratio,alt={Fig. 12}]{figs/fig12_H2_radiative_association.png}}
\textbf{Fig. 12.} Radiative association. Left: capture probability per
pass \(P_J(E)\) (\(J=0\)--4). Middle: \(\sigma(E)\) and \(k(E)\). Right:
\(k(E)\) and the thermal average \(k(T)\).

\subsection{5.7 Separate account of inter-atomic
gravity}\label{separate-account-of-inter-atomic-gravity}

The Newtonian gravity between the atoms,
\(\langle f_{vJ}|-Gm_H^2/R|f_{vJ}\rangle\) (point masses
\(m_H=m_p+m_e-E_b/c^2\)), is \(-1.244\times10^{-31}\) cm⁻¹ for (0,0) and
\(-3.97\times10^{-32}\) cm⁻¹ for (14,0); the line shifts are
\(10^{-37}\)--\(10^{-36}\) eV. This is 20 orders of magnitude below the
double-precision floor of the level energies (\(10^{-12}\) cm⁻¹), so
adding it to \(E\) changes nothing, and it is kept as a separate column
(Fig. 15). The gravitational cross terms electron--nucleus and
electron--electron (about \(10^{-3}\) and \(10^{-7}\) of the
nucleus--nucleus term; they require \(\sum\langle1/r_{ia}\rangle(R)\)
and \(\langle1/r_{12}\rangle(R)\)) and the 1PN terms
(\((v/c)^2\sim10^{-10}\)) are not included. In this system gravity
appears only as a shift of \(10^{-31}\) cm⁻¹ in the levels and as
gravitational-wave emission of \(10^{-35}\) eV.

\pandocbounded{\includegraphics[keepaspectratio,alt={Fig. 15}]{figs/fig15_H2_gravity_accounting.png}}
\textbf{Fig. 15.} Left: \(\langle-Gm_H^2/R\rangle\) per level. Right:
the ratio of the gravitational-wave to the E2 rate for each E2 line (the
ratio rises for lines with large \(\Delta v\) where
\(\langle\Theta\rangle\) is small by cancellation).

\begin{center}\rule{0.5\linewidth}{0.5pt}\end{center}

\section{6. Comparison of the four systems and
discussion}\label{comparison-of-the-four-systems-and-discussion}

\subsection{6.1 What remains in the
record}\label{what-remains-in-the-record}

{\def\LTcaptype{none} % do not increment counter
\begin{longtable}[]{@{}
  >{\raggedright\arraybackslash}p{(\linewidth - 6\tabcolsep) * \real{0.2500}}
  >{\raggedright\arraybackslash}p{(\linewidth - 6\tabcolsep) * \real{0.2500}}
  >{\raggedright\arraybackslash}p{(\linewidth - 6\tabcolsep) * \real{0.2500}}
  >{\raggedright\arraybackslash}p{(\linewidth - 6\tabcolsep) * \real{0.2500}}@{}}
\toprule\noalign{}
\begin{minipage}[b]{\linewidth}\raggedright
\end{minipage} & \begin{minipage}[b]{\linewidth}\raggedright
Hydrogen atom
\end{minipage} & \begin{minipage}[b]{\linewidth}\raggedright
Electron--electron, proton--proton
\end{minipage} & \begin{minipage}[b]{\linewidth}\raggedright
Hydrogen molecule
\end{minipage} \\
\midrule\noalign{}
\endhead
\bottomrule\noalign{}
\endlastfoot
bound levels & 50 (\(n\le5\)) & 0 & 301 \\
kinds of emission & E1, two-photon, M1, E2, GW & quadrupole
bremsstrahlung only (dipole vanishes identically) & E2, M1, GW (no
E1) \\
photons per cascade / collision & a few & \(5\times10^{-16}\) /
\(5\times10^{-23}\) & 2--7 \\
lowest transition vs \(kT\) at 2.7255 K & 21 cm \(=0.025\,kT\),
\(\bar n=39.5\) & --- & 587 cm⁻¹ \(=310\,kT\),
\(\bar n\sim10^{-135}\) \\
readout of layer 2 & spin temperature 2.72550 K & --- & none \\
readout of layer 1 & 0.09 m/s (ensemble), \(2.4\times10^{-10}\) m/s
(history of one atom) & distribution of deflection angles only & 0.002
m/s (line width \(10^{-22}\)) \\
readout of layer 3 & polarization of the initial \(m\), steady alignment
\(10^{-8}\) & --- & polarization of the initial \(m\), no relaxation \\
how gravity appears & levels \(10^{-38}\) eV, GW \(10^{-48}\) eV & GW/E2
\(10^{-43}\)--\(10^{-36}\) & levels \(10^{-31}\) cm⁻¹, GW \(10^{-35}\)
eV \\
\end{longtable}
}

The properties of the absorber enter the record in three different ways.
Temperature changes the generator (upward rates); direction and velocity
change only the readout of the record (Doppler and recoil); anisotropy
changes the state space, the generator and the readout. An isotropic,
infinitely distant, completely absorbing shell has no relational
quantity ``direction'', and neither translation nor recoil can be
defined. Direction is defined only as the assignment to the anonymous
equal-weight cells of the absorber, and is read from the angular
correlation of two-photon pairs and from the \(\cos\theta\) of the
dipole.

\subsection{6.2 Internal states cannot be
read}\label{internal-states-cannot-be-read}

The occupations of the 50 levels of the hydrogen atom, of the 301 levels
of H₂, and the intermediate orbit in scattering never appear in the
record. What appears is the sequence of emission and absorption events,
and two models that give the same record cannot be distinguished. The
bookkeeping of this paper (\(\tau\), the time evolution of occupations)
was computed for checking and is not observable. That assuming a real
orbit in the scattering of identical particles drops the Mott
interference term is one example.

\subsection{6.3 The structure of standard
theory}\label{the-structure-of-standard-theory}

In this set the interaction is used up at the moment the eigenvalue
problem is solved; during the run a constant generator \(R\) merely
moves probability. The exponential waiting-time distribution is
memoryless, and the system has no time of its own between events. The
line width is not an uncertainty of the level but the length of the
emission process, and it is Lorentzian if the envelope is exponential.
The two-photon pair is one event in which two quanta leave together;
only the sum frequency is sharp, each is continuous, and the
simultaneity of the pair carries the sharing of the origin to the
outside. These are the structure of the record that standard theory
gives for two-body systems, not claims of this paper. They are what a
future distribution-array operation must reproduce.

\subsection{6.4 The difference between the hydrogen atom and the
hydrogen
molecule}\label{the-difference-between-the-hydrogen-atom-and-the-hydrogen-molecule}

Running the same framework, the readout of the 2.7255 K absorber exists
only for the hydrogen atom. The hydrogen 21 cm line is 1/40 of \(kT\)
and equilibrates with the background, whereas the lowest transition of
H₂ is 310 times \(kT\) and the background neither excites nor relaxes
\(m\). The velocity recovery of layer 1 is more precise for H₂ because
the E2 lines are extremely narrow, \(10^{-22}\): 0.002 m/s with 200
molecules against 0.09 m/s with 200 hydrogen atoms.

\subsection{6.5 Time evolution of the field, the record, and the
physical counterpart of the
absorber}\label{time-evolution-of-the-field-the-record-and-the-physical-counterpart-of-the-absorber}

The generator of this paper has no field degrees of freedom. The static
field is folded into the eigenvalue problem, the radiation field into
the Einstein A coefficients, and a photon disappears into the absorber
at infinity at the moment of emission. What the time evolution of the
field leaves in the record is exactly the amount by which this paper
deviates from measurement: the 8170 MHz of 1S and the 1057.85 MHz of
2S--2P₁/₂ (Lamb shift), and the factors 0.9989 and 0.9944 of the
hyperfine structure and of 21 cm (\(g_e-2\)). Evolving the field
classically gives no stationary state and leads to collapse (the tenth
thought experiment of §1.2), so the record's ``the ground state does not
radiate'' requires both the time evolution of the field and a discrete
lock. This paper supplies the record that a generator with both must
reproduce.

The standpoint that takes the state of the field as the reality and the
standpoint that takes the record at infinity as the reality say the same
thing in different words as far as the radiative part is concerned
(§2.5). The two give the same predictions when the absorber is complete:
that is the condition under which the Wheeler--Feynman absorber theory
{[}58{]} agrees with the Maxwell radiation reaction, and under which
direct-action quantum electrodynamics {[}67{]} agrees with ordinary
quantum electrodynamics. If absorption is incomplete the two predict
differently. Partridge {[}59{]} looked for a directional dependence of
radiated power and found none, but Pegg {[}68{]} and others argued that
the null result was inevitable because the Earth itself acts as an
absorber, so that experiment does not decide whether absorption on a
cosmic scale is complete. The premise of this paper (a complete absorber
covering the whole sky) is not a fact established by measurement; it is
placed as the cosmological condition of Hogarth and Hoyle--Narlikar.

The physical counterpart of the absorber is the cosmological event
horizon, and its rest frame is the rest frame of the background
radiation. Hogarth {[}64{]} and Hoyle--Narlikar {[}65{]} argued that
whether radiation is completely absorbed is decided by the future
structure of the universe. The horizon of an accelerating universe
(radius \(c/H_0\approx4.4\) Gpc), seen from inside, is a complete
absorber covering all directions, and its area in Planck units,
\(4\pi(R_H/l_P)^2\approx7\times10^{122}\) (the de Sitter entropy
{[}66{]}), is the upper limit of the cell number \(N\) of §2.2. This
paper combines the incoming background (the past cone, 2.7255 K) and the
absorption (the future cone, with the horizon temperature
\(\hbar H_0/2\pi k\approx2\times10^{-30}\) K) into a heat bath of one
temperature. This is an approximation valid when the duration of the
experiment is shorter than the expansion time; in the runs to
\(t_{\max}=10^{17}\) s (3 billion years) the background temperature
changes by 20 \%, so strictly it is violated.

The only source of change of the internal phase is exchange with the
absorber. Recorded exchanges (emission and absorption) give the
frequencies themselves; unrecorded (virtual) exchanges give the shifts
of the frequencies (the Lamb shift). While there is no exchange with the
outside, the system has its eigenfrequencies, its own internal clock and
its spatial extent, but is stationary, and all that can be read is the
statistics of the next event. The phase of a superposition appears in
the time distribution of the next event as quantum beats, but the master
equation of this paper moves occupations only and does not produce them.

\subsection{6.6 Directions for further
work}\label{directions-for-further-work}

\begin{enumerate}
\def\labelenumi{\arabic{enumi}.}
\tightlist
\item
  \textbf{Correlations of the record.} The record of this paper stops at
  first order (intensities). The phase of the field appears as reality
  in correlations between records: \(g^{(1)}(\tau)\) (the line shape is
  properly its Fourier transform), \(g^{(2)}\) (antibunching of a single
  atom, \(g^{(2)}(0)=0\), and the time correlations of cascade photons),
  the polarization correlation of two-photon pairs (Bell {[}62{]}) and
  the interference of two sources. All can be computed in standard
  quantum optics and can be added to this paper as a second-order answer
  table.
\item
  \textbf{Traces of the vacuum field.} Bethe's nonrelativistic
  calculation (a mode sum of the radiation field with cutoff \(m_ec^2\))
  gives a 2S Lamb shift of 1040 MHz. The dipole matrix elements and
  Bethe logarithms it needs are already in the present set.
\item
  \textbf{Separation of past and future.} Hold the temperature of the
  incoming background and the future absorption as separate surfaces and
  include the change of temperature by expansion.
\item
  \textbf{The generator.} A generator (distribution-array operation)
  that produces both radiation and the lock from one exchange must
  reproduce the record of this paper and its correlations. That is the
  goal of this research series.
\end{enumerate}

\begin{center}\rule{0.5\linewidth}{0.5pt}\end{center}

\section{7. Limitations and
non-claims}\label{limitations-and-non-claims}

This paper does not claim the following.

\begin{enumerate}
\def\labelenumi{\arabic{enumi}.}
\tightlist
\item
  \textbf{That the Einstein--Maxwell equations were implemented.} The
  field is not evolved in time. The levels are bound-state formulas,
  gravity is 1PN expectation values, radiation is Einstein A
  coefficients.
\item
  \textbf{That the dynamics was derived from a distribution-array
  operation of a universal interaction.} This paper only supplies the
  target of comparison.
\item
  \textbf{The derivation of new physical laws or constants.} All inputs
  are known values (§2.1).
\item
  \textbf{Hydrogen levels including QED and nuclear size.} Apart from
  the 0.12 MHz of 1S, every mismatch equals the size of the omitted QED
  and nuclear-size terms (§3.3). Gravity beyond 2PN is not included.
\item
  \textbf{H₂ levels including nonadiabatic, relativistic and QED
  corrections.} The levels differ from the fully corrected values by
  −0.3 to +4.6 cm⁻¹ and (14,4) is not bound. \(\Theta(R)\) is the 1998
  value, and the A coefficients of strong lines differ by up to 6 \%.
\item
  \textbf{The triplet channel and electronically excited states of H₂.}
  b³Σu⁺ (repulsive, 3/4 of electron spin) and the Lyman and Werner bands
  are not included.
\item
  \textbf{A quantum theory of electron--electron bremsstrahlung.} At
  \(\eta\approx3.5\) the classical spectrum is a guide; no Gaunt factor
  is included.
\item
  \textbf{A complete account of inter-atomic gravity.} Only the
  point-mass Newtonian term; the electron cross terms and 1PN are not
  included (§5.7).
\item
  \textbf{The actual axis of the absorber quadrupole.}
  \(a_2=4\times10^{-6}\) is a guide value from the COBE \(Q_{\rm rms}\)
  and is placed on the same axis as the dipole.
\item
  \textbf{The widths of the radiative-association resonances.} The box
  discretization does not resolve the widths of shape resonances. The
  (14,4) resonance disappears when the nonadiabatic correction is
  included.
\item
  \textbf{A physical derivation of the cell number \(N=10^{60}\).} It
  was placed as a value sufficiently above the saturation lower limit
  \(1.5\times10^{34}\) and below the area-law upper limit
  \(4\pi(R_H/l_P)^2\approx7\times10^{122}\). The value does not appear
  in the record.
\item
  \textbf{The identification of the absorber with the cosmological
  horizon, and the premise of complete absorption.} The identification
  of §6.5 is an interpretation; the computation was done with a static
  absorber of one temperature. Complete absorption is a condition
  placed, not a fact established by measurement. The separation of the
  past background from the future absorption and the change of
  temperature by expansion are not included.
\item
  \textbf{Correlations of the record.} The record of this paper is
  intensities (first order) only; \(g^{(1)}\), \(g^{(2)}\), the
  polarization correlation of two-photon pairs and quantum beats were
  not computed.
\end{enumerate}

\begin{center}\rule{0.5\linewidth}{0.5pt}\end{center}

\section{8. Reproducibility}\label{reproducibility}

All programs, data and figures are in the GitHub repository
ai-chat-logs-open under
\texttt{匿名頂点状態生成幾何-匿名内部観測者から不変な関係量の体系/Grok移行実験/};
each folder regenerates its set with \texttt{run\_all.sh} and carries a
\texttt{SHA256SUMS}. The same files are included in the Zenodo record.

{\def\LTcaptype{none} % do not increment counter
\begin{longtable}[]{@{}
  >{\raggedright\arraybackslash}p{(\linewidth - 4\tabcolsep) * \real{0.3226}}
  >{\raggedright\arraybackslash}p{(\linewidth - 4\tabcolsep) * \real{0.4839}}
  >{\raggedright\arraybackslash}p{(\linewidth - 4\tabcolsep) * \real{0.1935}}@{}}
\toprule\noalign{}
\begin{minipage}[b]{\linewidth}\raggedright
Folder
\end{minipage} & \begin{minipage}[b]{\linewidth}\raggedright
Contents
\end{minipage} & \begin{minipage}[b]{\linewidth}\raggedright
Reproduction
\end{minipage} \\
\midrule\noalign{}
\endhead
\bottomrule\noalign{}
\endlastfoot
\texttt{exchange\_rel\_20261005/} & the fixed hydrogen set:
\texttt{exchange\_cascade.py} (levels, A coefficients, master equation),
three checks (\texttt{coulomb\_levels\_check.py},
\texttt{gravity\_levels\_check.py}, \texttt{rates\_check.py}), the
tables \texttt{structure\_tables.py}, the structure document
\texttt{プログラム構造\_ja\_20261006.md}, the gravity survey
\texttt{重力項目\_調査と修正\_ja\_20261005.md}, the diff from the
original \texttt{CHANGES\_vs\_original.diff} & \texttt{run\_all.sh}, SHA
16 files \\
\texttt{exchange\_general\_20261006/} & the sign-input numerical solver
\texttt{two\_body\_general.py} (ep / ee / pp), the control
\texttt{control\_general\_vs\_closed\_form.py}, scattering
\texttt{scattering\_same\_sign.py}, the three absorber layers
\texttt{absorber\_temperature.py},
\texttt{absorber\_direction\_recoil.py}, \texttt{absorber\_alignment.py}
& \texttt{run\_all.sh}, SHA 19 files \\
\texttt{exchange\_engine\_20261006/} & the common engine
\texttt{engine.py}, the hydrogen component with its regression test, the
H₂ component \texttt{system\_h2.py}, validation
\texttt{validation\_h2.py}, runs \texttt{run\_h2.py}, radiative
association \texttt{radiative\_association\_h2.py}, literature data
\texttt{h2\_data/} (sources and SHA of the downloaded files in
\texttt{SOURCES.md}), JSON for the figures & \texttt{run\_all.sh} (about
5 min), SHA 29 files \\
\texttt{figures\_paper\_20261006/} & the 15 figures of this paper,
\texttt{make\_figures.py}, the plotted numbers \texttt{data/*.json} &
\texttt{run\_all.sh}, SHA 29 files \\
\end{longtable}
}

Main checks: the regression of the engine version against the fixed
hydrogen set (difference of the emission 13.054539268 eV
\(-3.6\times10^{-15}\); differences of generator, dwell times and final
state 0), the control of the fixed set against the numerical solver
(§3.6), the comparison of H₂ with literature (§5.3), the box
independence and the round-trip time of the radiative association
(§5.6), and energy conservation (total emission
\(=E_{\rm init}-E_{\rm final}\), difference \(\le10^{-14}\) eV).

Random sequences: \texttt{default\_rng(7)} for the layer-1 ensemble,
\texttt{default\_rng(20261006)} for the one-atom history,
\texttt{default\_rng(11)} for layer 1 of H₂. Environment: Python 3.9,
numpy 2.0.2, scipy, matplotlib. On macOS with Accelerate, \texttt{@} on
matrices of 50×50 or larger emits a false warning, so \texttt{np.dot} is
used. The distribution site of H2SPECTRE 7.4 has an expired certificate,
and the archive contains duplicated hard links, so only the regular
files were extracted with Python's tarfile. Running heavy computations
in parallel makes Accelerate contend for all cores and slows them down
by more than a factor of ten.

Every number in this paper is transcribed as it stands from the values
written in \texttt{audit*.txt}, \texttt{*.md} and \texttt{*.json} of
these folders.

\begin{center}\rule{0.5\linewidth}{0.5pt}\end{center}

\section{9. Record of AI involvement}\label{record-of-ai-involvement}

\begin{itemize}
\tightlist
\item
  \textbf{Grok} (2026-10-05): the first program of the multistep
  emission of the hydrogen atom (a series cascade). The present set
  started from a copy of it and replaced the Coulomb levels with
  Dirac--Coulomb plus recoil, gravity with the 1PN of general
  relativity, and the dynamics with the master equation of Einstein A
  coefficients. Grok's rates (\(\propto\Delta E\)), spin terms, rate
  feedback and the \(n\to n-1\) constraint were all removed. The diff is
  \texttt{CHANGES\_vs\_original.diff}.
\item
  \textbf{ChatGPT} (2026-09-30 to 10-01): the Einstein--Maxwell complete
  specification (v1.0--v3.1) that preceded this paper. Not used here.
\item
  \textbf{Claude} (2026-10-05 to 07): the above corrections and
  implementation, the numerical solver, scattering, the three layers,
  the engine, H₂, radiative association, the figures and the drafting of
  this paper. Design and judgement are the author's.
\end{itemize}

\begin{center}\rule{0.5\linewidth}{0.5pt}\end{center}

\section{10. Conclusion}\label{conclusion}

A program was built that computes the emission record standard theory
gives for two-body systems, holding neither a coordinate grid nor a
metric of a background spacetime and referenced only to the rest frame
of an absorber at infinity, and it was run on four systems: the hydrogen
atom, electron--electron, proton--proton and the hydrogen molecule. The
hydrogen atom agrees with external measurements to the precision that
excludes QED and nuclear size, and the hydrogen molecule agrees with the
literature A coefficients to \(10^{-3}\). Same-sign pairs have no bound
levels, and only the Mott cross section and quadrupole bremsstrahlung
remain in the record. The temperature, direction and velocity, and
anisotropy of the absorber enter the record each in a different way; in
the 21 cm line of the hydrogen atom the temperature of the absorber can
be read, in H₂ it cannot. In every system gravity appears only at
\(10^{-31}\)--\(10^{-38}\) of the levels and as
\(10^{-35}\)--\(10^{-48}\) eV of gravitational waves.

This record is the reference that a future distribution-array operation
must reproduce.

\begin{center}\rule{0.5\linewidth}{0.5pt}\end{center}

\section{Appendix A. States and edges of the hydrogen
set}\label{appendix-a.-states-and-edges-of-the-hydrogen-set}

50 levels (\(n\le5\), \((n,\ell,s_e,s_p)\)), initial 5p₃/₂ F=1. Edges:
E1 320, E2 376, GW 376, M1 59, two-photon 2, in total 1133. The full
tables of levels, edges, emission and widths are in
\texttt{exchange\_rel\_20261005/structure\_tables.md} (Tables A--F), and
the correspondence between formulas and code is in
\texttt{プログラム構造\_ja\_20261006.md}.

\section{Appendix B. Formulas of the absorber
layers}\label{appendix-b.-formulas-of-the-absorber-layers}

\begin{itemize}
\tightlist
\item
  Doppler:
  \(\delta=\omega_{\rm lab}/\omega'-1=\dfrac{\beta n_\parallel\gamma-(\gamma-1)}{\gamma(1-\beta n_\parallel)}\)
  (without cancellation).
\item
  Recoil: exact two-body kinematics in the rest frame, emission
  \(\hbar\omega'=\Delta E(1-\Delta E/2Mc^2)\), absorption
  \(\Delta E(1+\Delta E/2M_fc^2)\). \(\vec v\) is updated from the
  momentum \(\vec P\to\vec P\mp(\hbar\omega_{\rm lab}/c)\hat n\) and the
  energy \(E\to E\mp\hbar\omega_{\rm lab}\).
\item
  Angular patterns \(P_q(\theta)\) of layer 3: for \(k=1\), \(q=0\):
  \((3/8\pi)\sin^2\theta\), \(q=\pm1\): \((3/16\pi)(1+\cos^2\theta)\).
  For \(k=2\), \(q=0\): \((15/8\pi)\sin^2\theta\cos^2\theta\),
  \(|q|=1\): \((5/16\pi)(1-3\cos^2\theta+4\cos^4\theta)\), \(|q|=2\):
  \((5/16\pi)(1-\cos^4\theta)\).
\item
  Occupation number by type
  \(\langle\bar n\rangle_q=\int\bar n(\omega,T(\hat n))P_q\,d\Omega\);
  steady sublevels
  \(p_m\propto\langle\bar n\rangle_q/(1+\langle\bar n\rangle_q)\).
\end{itemize}

\section{Appendix C. List of figures}\label{appendix-c.-list-of-figures}

{\def\LTcaptype{none} % do not increment counter
\begin{longtable}[]{@{}
  >{\raggedright\arraybackslash}p{(\linewidth - 2\tabcolsep) * \real{0.0312}}
  >{\raggedright\arraybackslash}p{(\linewidth - 2\tabcolsep) * \real{0.9688}}@{}}
\toprule\noalign{}
\begin{minipage}[b]{\linewidth}\raggedright
Fig.
\end{minipage} & \begin{minipage}[b]{\linewidth}\raggedright
File
\end{minipage} \\
\midrule\noalign{}
\endhead
\bottomrule\noalign{}
\endlastfoot
1 & \texttt{figures/fig01\_H\_levels\_cascade.png} \\
2 & \texttt{figures/fig02\_H\_spectrum\_widths.png} \\
3 & \texttt{figures/fig03\_H\_time\_evolution.png} \\
4 & \texttt{figures/fig04\_H\_layer2\_temperature.png} \\
5 & \texttt{figures/fig05\_H\_layer1\_doppler\_recoil.png} \\
6 & \texttt{figures/fig06\_H\_layer3\_alignment.png} \\
7 & \texttt{figures/fig07\_ee\_pp\_mott.png} \\
8 & \texttt{figures/fig08\_ee\_pp\_bremsstrahlung.png} \\
9 & \texttt{figures/fig09\_H2\_levels\_cascades.png} \\
10 & \texttt{figures/fig10\_H2\_cascade\_spectra.png} \\
11 & \texttt{figures/fig11\_H2\_validation.png} \\
12 & \texttt{figures/fig12\_H2\_radiative\_association.png} \\
13 & \texttt{figures/fig13\_H2\_layer1\_doppler.png} \\
14 & \texttt{figures/fig14\_H2\_layer3\_polarization.png} \\
15 & \texttt{figures/fig15\_H2\_gravity\_accounting.png} \\
16 & \texttt{figures/fig16\_lightcone\_absorber.png} (schematic,
\texttt{figures/make\_fig16\_lightcone.py}) \\
\end{longtable}
}

Figures 1--15 are generated by
\texttt{Grok移行実験/figures\_paper\_20261006/make\_figures.py}, and the
plotted numbers are in \texttt{data/*.json} of that folder. Figure 16 is
a schematic drawn by \texttt{figures/make\_fig16\_lightcone.py} in this
folder and has no data behind it.

\begin{center}\rule{0.5\linewidth}{0.5pt}\end{center}

\section{References}\label{references}

\subsection{This research series}\label{this-research-series}

{[}S1{]} N. Kihara, \textbf{Designing a Research Project to Explore the
Emergence of Spacetime, Particles, and Interactions from States and
Relations: A Framework of Preliminary Design and Feasibility
Verification for Foundational-Physics Model Exploration under
Uncertainty} (Paper 0, the design document). Concept DOI:
10.5281/zenodo.22851944.

{[}S2{]} N. Kihara, \textbf{The Sixth Thought Experiment: Can Two Kinds
of Long-Range Interaction Be Internalized into a Single Closed State
Map? - Reconstructing Gravity, Coulomb, GW-Quadrupole and EM-Dipole
Radiation of a Charged Quasi-Circular Binary with Eight Persistent
States, and Verifying the Transferability of the Same Transition over
Five Conditions}, v1.1 (2026-09-26). Concept DOI:
10.5281/zenodo.22974631; Version DOI: 10.5281/zenodo.22974632.

{[}S3{]} N. Kihara, \textbf{Supplement to the Sixth Thought Experiment:
Does the State Generation Change When the Normalization Reference Is
Changed from Mass to a Charge Unit? - A Restricted Renormalization from
G=M=c=1 to G=q₀=c=1, and an Exact-Match Verification of the Initial
States of the Five Stored Conditions}, v1.0 (2026-09-27). Concept DOI:
10.5281/zenodo.22985299; Version DOI: 10.5281/zenodo.22985300.

{[}S4{]} N. Kihara, \textbf{Composition of Central Projection and the
Closed Form of the Composite Curvature Radius: An Algebraic Formulation
of High-Dimensional Reduction via One Central Projection and Commutative
Cuts on the Sphere}, v1 (2026-05-07). Concept DOI:
10.5281/zenodo.20060728; Version DOI: 10.5281/zenodo.20060729.

{[}S5{]} N. Kihara, \textbf{Radial Projection: Definition and Relation
to Central Projection}, v3.3 (2026-06-06). Concept DOI:
10.5281/zenodo.20462569; Version DOI: 10.5281/zenodo.20567347.

Records of thought experiments 7--10 and of the Einstein--Maxwell
discrete variation (unpublished): the folders
\texttt{第七思考実験\_自己相互作用\_aa\_bb\_20260927/},
\texttt{第八思考実験\_aa\_ab\_bb\_交差結合仮定探索\_20260927/},
\texttt{第九思考実験\_倍音干渉による有限履歴読出し\_20260929/},
\texttt{第十思考実験\_EM自己無撞着有限履歴\_LW\_LAD\_20260929/} of the
same repository.

\subsection{External literature}\label{external-literature}

{[}1{]} P. J. Mohr, D. B. Newell, B. N. Taylor, E. Tiesinga, CODATA
recommended values of the fundamental physical constants: 2022,
Rev.~Mod. Phys. 97, 025002 (2025).

{[}2{]} E. Tiesinga, P. J. Mohr, D. B. Newell, B. N. Taylor, CODATA
recommended values of the fundamental physical constants: 2018,
Rev.~Mod. Phys. 93, 025010 (2021).

{[}3{]} P. J. Mohr, B. N. Taylor, D. B. Newell, CODATA recommended
values of the fundamental physical constants: 2010, Rev.~Mod. Phys. 84,
1527 (2012).

{[}4{]} P. A. M. Dirac, The quantum theory of the electron, Proc. R.
Soc. London A 117, 610 (1928).

{[}5{]} A. Sommerfeld, Zur Quantentheorie der Spektrallinien, Ann. Phys.
51, 1 (1916).

{[}6{]} W. A. Barker, F. N. Glover, Reduction of relativistic
two-particle wave equations to approximate forms. III, Phys. Rev.~99,
317 (1955).

{[}7{]} G. Breit, Possible effects of nuclear spin on X-ray terms, Phys.
Rev.~35, 1447 (1930).

{[}8{]} E. E. Salpeter, Mass corrections to the fine structure of
hydrogen-like atoms, Phys. Rev.~87, 328 (1952).

{[}9{]} G. W. Erickson, Energy levels of one-electron atoms, J. Phys.
Chem. Ref. Data 6, 831 (1977).

{[}10{]} J. R. Sapirstein, D. R. Yennie, Theory of hydrogenic bound
states, in Quantum Electrodynamics, ed.~T. Kinoshita (World Scientific,
Singapore, 1990), p.~560.

{[}11{]} K. Pachucki, H. Grotch, Pure recoil corrections to hydrogen
energy levels, Phys. Rev.~A 51, 1854 (1995).

{[}12{]} M. I. Eides, H. Grotch, Recoil corrections of order (Zα)⁶(m/M)m
to the hydrogen energy levels, Phys. Rev.~A 55, 3351 (1997).

{[}13{]} E. A. Golosov, I. B. Khriplovich, A. I. Milstein, A. S.
Yelkhovsky, Order α⁴(m/M)R∞ corrections to hydrogen P levels, JETP 80,
208 (1995).

{[}14{]} U. Jentschura, K. Pachucki, Higher-order binding corrections to
the Lamb shift of 2P states, Phys. Rev.~A 54, 1853 (1996).

{[}15{]} K. Pachucki, S. G. Karshenboim, Higher order recoil corrections
to energy levels of two-body systems, Phys. Rev.~A 60, 2792 (1999).

{[}16{]} K. Melnikov, A. Yelkhovsky, O(mα⁷ln²α) corrections to
positronium energy levels, Phys. Lett. B 458, 143 (1999).

{[}17{]} G. W. F. Drake, R. A. Swainson, Bethe logarithms for hydrogen
up to n = 20, and approximations for two-electron atoms, Phys. Rev.~A
41, 1243 (1990).

{[}18{]} U. D. Jentschura, P. J. Mohr, Calculation of hydrogenic Bethe
logarithms for Rydberg states, Phys. Rev.~A 72, 012110 (2005).

{[}19{]} M. I. Eides, H. Grotch, V. A. Shelyuto, Theory of light
hydrogenlike atoms, Phys. Rep.~342, 63 (2001).

{[}20{]} E. Fermi, Über die magnetischen Momente der Atomkerne, Z. Phys.
60, 320 (1930).

{[}21{]} A. Einstein, L. Infeld, B. Hoffmann, The gravitational
equations and the problem of motion, Ann. Math. 39, 65 (1938).

{[}22{]} B. M. Barker, R. F. O'Connell, Gravitational two-body problem
with arbitrary masses, spins, and quadrupole moments, Phys. Rev.~D 12,
329 (1975).

{[}23{]} Yu. N. Obukhov, Spin, gravity, and inertia, Phys. Rev.~Lett.
86, 192 (2001).

{[}24{]} L. L. Foldy, S. A. Wouthuysen, On the Dirac theory of spin 1/2
particles and its non-relativistic limit, Phys. Rev.~78, 29 (1950).

{[}25{]} M. Khalil, N. Sennett, J. Steinhoff, J. Vines, A. Buonanno,
Hairy binary black holes in Einstein-Maxwell-dilaton theory and their
effective-one-body description, Phys. Rev.~D 98, 104010 (2018).

{[}26{]} P. K. Gupta, Binary dynamics from Einstein-Maxwell theory at
second post-Newtonian order using effective field theory, Phys. Rev.~D
112, 104047 (2025); arXiv:2205.11591.

{[}27{]} L. D. Landau, E. M. Lifshitz, The Classical Theory of Fields,
4th ed.~(Butterworth-Heinemann, Oxford, 1975), §71, §110.

{[}28{]} A. Einstein, Zur Quantentheorie der Strahlung, Phys. Z. 18, 121
(1917).

{[}29{]} I. I. Sobelman, Atomic Spectra and Radiative Transitions, 2nd
ed.~(Springer, Berlin, 1992).

{[}30{]} A. R. Edmonds, Angular Momentum in Quantum Mechanics (Princeton
University Press, Princeton, 1957).

{[}31{]} G. Racah, Theory of complex spectra. II, Phys. Rev.~62, 438
(1942).

{[}32{]} G. Breit, E. Teller, Metastability of hydrogen and helium
levels, Astrophys. J. 91, 215 (1940).

{[}33{]} S. P. Goldman, Generalized Laguerre representation: Application
to relativistic two-photon decay rates, Phys. Rev.~A 40, 1185 (1989).

{[}34{]} W. R. Johnson, Radiative decay rates of metastable one-electron
atoms, Phys. Rev.~Lett. 29, 1123 (1972).

{[}35{]} A. Kramida, Yu. Ralchenko, J. Reader, NIST ASD Team, NIST
Atomic Spectra Database (ver. 5.12), https://physics.nist.gov/asd
(accessed October 2026).

{[}36{]} W. L. Wiese, J. R. Fuhr, Accurate atomic transition
probabilities for hydrogen, helium, and lithium, J. Phys. Chem. Ref.
Data 38, 565 (2009).

{[}37{]} A. E. Kramida, A critical compilation of experimental data on
spectral lines and energy levels of hydrogen, deuterium, and tritium,
At. Data Nucl. Data Tables 96, 586 (2010).

{[}38{]} E. W. Hagley, F. M. Pipkin, Separated oscillatory field
measurement of hydrogen 2S₁/₂--2P₃/₂ fine structure interval, Phys.
Rev.~Lett. 72, 1172 (1994).

{[}39{]} C. G. Parthey et al., Improved measurement of the hydrogen
1S--2S transition frequency, Phys. Rev.~Lett. 107, 203001 (2011).

{[}40{]} L. Essen, R. W. Donaldson, M. J. Bangham, E. G. Hope, Frequency
of the hydrogen maser, Nature 229, 110 (1971).

{[}41{]} D. J. Fixsen, The temperature of the cosmic microwave
background, Astrophys. J. 707, 916 (2009).

{[}42{]} Planck Collaboration, Planck 2018 results. I. Overview and the
cosmological legacy of Planck, Astron. Astrophys. 641, A1 (2020).

{[}43{]} C. L. Bennett et al., Four-year COBE DMR cosmic microwave
background observations: Maps and basic results, Astrophys. J. 464, L1
(1996).

{[}44{]} N. F. Mott, The collision between two electrons, Proc. R. Soc.
London A 126, 259 (1930).

{[}45{]} E. Rutherford, The scattering of α and β particles by matter
and the structure of the atom, Philos. Mag. 21, 669 (1911).

{[}46{]} K. Pachucki, Born-Oppenheimer potential for H₂, Phys. Rev.~A
82, 032509 (2010).

{[}47{]} J. Komasa, M. Puchalski, P. Czachorowski, G. Łach, K. Pachucki,
Rovibrational energy levels of the hydrogen molecule through
nonadiabatic perturbation theory, Phys. Rev.~A 100, 032519 (2019);
H2SPECTRE ver. 7.4 (2022), https://qcg.home.amu.edu.pl/H2Spectre.html.

{[}48{]} K. Pachucki, J. Komasa, Accurate adiabatic correction in the
hydrogen molecule, J. Chem. Phys. 141, 224103 (2014).

{[}49{]} L. Wolniewicz, I. Simbotin, A. Dalgarno, Quadrupole transition
probabilities for the excited rovibrational states of H₂, Astrophys. J.
Suppl. Ser. 115, 293 (1998).

{[}50{]} K. Pachucki, J. Komasa, Magnetic dipole transitions in the
hydrogen molecule, Phys. Rev.~A 83, 032501 (2011).

{[}51{]} E. Roueff, H. Abgrall, P. Czachorowski, K. Pachucki, M.
Puchalski, J. Komasa, The full infrared spectrum of molecular hydrogen,
Astron. Astrophys. 630, A58 (2019); CDS J/A+A/630/A58.

{[}52{]} K. Pachucki, J. Komasa, Nonadiabatic rotational states of the
hydrogen molecule, Phys. Chem. Chem. Phys. 20, 247 (2018).

{[}53{]} K. T. Tang, J. P. Toennies, An improved simple model for the
van der Waals potential based on universal damping functions for the
dispersion coefficients, J. Chem. Phys. 80, 3726 (1984).

{[}54{]} Z.-C. Yan, J. F. Babb, A. Dalgarno, G. W. F. Drake, Variational
calculations of dispersion coefficients for interactions among H, He,
and Li atoms, Phys. Rev.~A 54, 2824 (1996).

{[}55{]} D. T. Colbert, W. H. Miller, A novel discrete variable
representation for quantum mechanical reactive scattering via the
S-matrix Kohn method, J. Chem. Phys. 96, 1982 (1992).

{[}56{]} W. Kołos, L. Wolniewicz, Potential-energy curves for the X¹Σg⁺,
b³Σu⁺, and C¹Πu states of the hydrogen molecule, J. Chem. Phys. 43, 2429
(1965).

{[}57{]} E. P. Wigner, On the behavior of cross sections near
thresholds, Phys. Rev.~73, 1002 (1948).

{[}58{]} J. A. Wheeler, R. P. Feynman, Interaction with the absorber as
the mechanism of radiation, Rev.~Mod. Phys. 17, 157 (1945).

{[}59{]} R. B. Partridge, Absorber theory of radiation and the future of
the universe, Nature 244, 263 (1973).

{[}60{]} R. Arnowitt, S. Deser, C. W. Misner, The dynamics of general
relativity, in Gravitation: An Introduction to Current Research, ed.~L.
Witten (Wiley, New York, 1962).

{[}61{]} H. Bondi, M. G. J. van der Burg, A. W. K. Metzner,
Gravitational waves in general relativity. VII. Waves from axi-symmetric
isolated systems, Proc. R. Soc. London A 269, 21 (1962).

{[}62{]} W. Perrie, A. J. Duncan, H. J. Beyer, H. Kleinpoppen,
Polarization correlation of the two photons emitted by metastable atomic
deuterium: A test of Bell's inequality, Phys. Rev.~Lett. 54, 1790
(1985).

{[}63{]} W. Heisenberg, Die „beobachtbaren Größen'' in der Theorie der
Elementarteilchen, Z. Phys. 120, 513 (1943).

{[}64{]} J. E. Hogarth, Cosmological considerations of the absorber
theory of radiation, Proc. R. Soc. London A 267, 365 (1962).

{[}65{]} F. Hoyle, J. V. Narlikar, Time symmetric electrodynamics and
the arrow of time in cosmology, Proc. R. Soc. London A 277, 1 (1964).

{[}66{]} G. W. Gibbons, S. W. Hawking, Cosmological event horizons,
thermodynamics, and particle creation, Phys. Rev.~D 15, 2738 (1977).

{[}67{]} P. C. W. Davies, Extension of Wheeler-Feynman quantum theory to
the relativistic domain. I. Scattering processes, J. Phys. A 4, 836
(1971); II. Emission processes, J. Phys. A 5, 1025 (1972).

{[}68{]} D. T. Pegg, Absorber theory of radiation, Rep.~Prog. Phys. 38,
1339 (1975).

\begin{center}\rule{0.5\linewidth}{0.5pt}\end{center}

\section{Revision history}\label{revision-history}

\begin{itemize}
\tightlist
\item
  v1.0 (2026-10-07): first version. Translated from the Japanese
  original of the same date.
\end{itemize}

\end{document}
